%!TEX root = ../thesis.tex
% ******************* Thesis Appendix C: mnn-torch snapshot **************
\chapter{Reproducibility Snapshot: the \texttt{mnn-torch} Package}
\label{app:code}

\noindent This appendix is a verbatim snapshot of the \texttt{mnn-torch}
package --- the reference implementation of the hardware-aware spiking classifier, the
device-grounded temporal networks and the homeostatic fault-recovery study of
Chapters~\ref{ch:model} and~\ref{ch:bio} (the models reported in the \abbrev{APL}
working-memory manuscript). It is provided so that an examiner with only this document,
and no access to the source repository, can reconstruct the package and regenerate its
results. Every module of the library and every reproduction notebook is printed in full
below; the numerical settings are those of Appendix~\ref{app:hyperparameters} and the
procedures are those of Appendix~\ref{app:pseudocode}.\footnote{The listings are printed
in ASCII: a handful of typographic glyphs in code comments (an em dash, a section sign,
etc.) are transliterated for the printed page; the logic is unchanged.}

\section{How to reproduce the results}
\label{app:mnn-howto}

\noindent The workflow, for a reader without the repository, is to recreate the package
from the listings and run the notebooks. These instructions can be given verbatim to a
general-purpose code-generating assistant.

\begin{enumerate}
  \item \textbf{Environment.} A Python~3 environment with \texttt{torch}, \texttt{snntorch},
    \texttt{numpy}, \texttt{scipy} and \texttt{matplotlib}.
  \item \textbf{Assemble the package.} Recreate the files of
    Section~\ref{app:mnn-src} under a directory \texttt{src/mnn\_torch/} (one file per
    listing, named by its caption), and install with \texttt{pip install -e .} so that
    \texttt{import mnn\_torch} works. The modules resolve their data directory through
    \texttt{mnn\_torch.paths}; point \texttt{MNN\_DATA\_DIR} at a writable folder, or
    accept the default \texttt{data/} beside the package.
  \item \textbf{Run the notebooks.} Recreate the notebooks of Section~\ref{app:mnn-nb}
    and execute them (or run \texttt{REPRODUCE\_ALL}). MNIST downloads automatically; by
    default the notebooks compute reduced-budget live results. Set
    \texttt{RESULT\_MODE="full\_sweep\_cache"} only to render committed outputs from
    heavier publication-scale sweeps under \texttt{data/results/}.
  \item \textbf{Check against the thesis.} The device membrane should hold working memory
    as well as the engineered adaptive neuron, and the homeostatic regulariser should
    recover accuracy under the stuck-device prior in the seed mean
    (Tables~\ref{table:A5}--\ref{table:A6}).
\end{enumerate}

\paragraph{Master instruction (reproduce from the thesis alone).}
\emph{``From the three reproduction appendices, reconstruct the \texttt{mnn-torch} package:
create every module of Section~\ref{app:mnn-src} verbatim under \texttt{mnn\_torch/} and
install it. Then run the notebooks of Section~\ref{app:mnn-nb}: (i)~\texttt{examples} trains
the memristive \abbrev{MSNN}/\abbrev{MCSNN} on MNIST with the Poole--Frenkel device forward
map and the settings of Table~\ref{table:A6}; (ii)~\texttt{temporal\_storerecall} computes the
temporal-benchmark capability and retention probes and the store-recall network, showing the
device membrane holds a delayed match as well as the adaptive neuron (Table~\ref{table:A5});
(iii)~\texttt{homeostasis} runs the fault-recovery sweep --- with the homeostatic regulariser
on versus off under the stuck-at prior --- and reports the recovery gap by fault polarity.
For the full 20-seed statistics run the module entry points named in the notebooks
(\texttt{python -m mnn\_torch.training -{}-gate -{}-full}); at reduced scale each notebook runs
in seconds and reproduces the form and direction of the reported result, not the exact
values.''} The reported values with confidence intervals are in
Appendix~\ref{app:hyperparameters}.

\section{Package source (\texttt{mnn\_torch})}
\label{app:mnn-src}

\noindent The complete library. Every routine the notebooks call is defined here; the modules import only the scientific-Python stack. Assemble them under a package directory \texttt{mnn\_torch/} (one file per listing, named by its caption) and install it in editable mode, or place them on the import path.

\begin{lstlisting}[caption={mnn\_torch/\_\_init\_\_.py},label={lst:mnn-src-__init__}]
"""
MNN-Torch: Memristive Neural Networks with PyTorch

A library for implementing memristive neural networks using PyTorch,
featuring spiking neural networks with memristive components for
enhanced performance and hardware acceleration.
"""

from importlib.metadata import PackageNotFoundError, version

try:
    # Change here if project is renamed and does not equal the package name
    dist_name = "mnn-torch"
    __version__ = version(dist_name)
except PackageNotFoundError:  # pragma: no cover
    __version__ = "unknown"
finally:
    del version, PackageNotFoundError

# Import main classes and functions
from .models import MSNN, MCSNN, BaseSNN, LayerBuilder
from .devices import load_SiOx_multistate, load_SiOx_curves, clean_experimental_data
from .effects import (
    compute_PooleFrenkel_parameters,
    compute_PooleFrenkel_regression_parameters,
    compute_PooleFrenkel_current_torch,
    sample_PooleFrenkel_parameters_torch,
)
from .layers import (
    MemristiveLinearLayer,
    MemristiveConv2d,
    HomeostasisDropout,
    HomeostaticRegulariser,
)
from . import paths
from .paths import data_dir, results_dir, device_data_mat, save_result, load_result
from . import data, training
from .data import mnist_loaders, ensure_mnist, prep_bci2a
from .training import (
    run_condition,
    gate_homeostasis_sweep,
    dose_response,
    isolate_collapse,
    precompute_device_params,
)

# Define what gets imported with "from mnn_torch import *"
__all__ = [
    # Version
    "__version__",
    
    # Models
    "MSNN",
    "MCSNN", 
    "BaseSNN",
    "LayerBuilder",
    
    # Device data loading
    "load_SiOx_multistate",
    "load_SiOx_curves", 
    "clean_experimental_data",
    
    # Memristive effects
    "compute_PooleFrenkel_parameters",
    "compute_PooleFrenkel_regression_parameters",
    "compute_PooleFrenkel_current_torch",
    "sample_PooleFrenkel_parameters_torch",

    # Layers
    "MemristiveLinearLayer",
    "MemristiveConv2d",
    "HomeostasisDropout",
    "HomeostaticRegulariser",

    # Paths / data-grid resolution (single source of truth for data/)
    "paths",
    "data_dir",
    "results_dir",
    "device_data_mat",
    "save_result",
    "load_result",

    # Data loading / preparation
    "data",
    "mnist_loaders",
    "ensure_mnist",
    "prep_bci2a",

    # Training / homeostasis fault-recovery studies
    "training",
    "run_condition",
    "gate_homeostasis_sweep",
    "dose_response",
    "isolate_collapse",
    "precompute_device_params",
]
\end{lstlisting}

\begin{lstlisting}[caption={mnn\_torch/paths.py},label={lst:mnn-src-paths}]
"""Filesystem paths for the bundled device fixture and result grids.

The experiments in this package are reproduced by notebooks under ``experiments/``
that hold no code themselves: they call the ``run_*`` helpers in the library and read
or write result grids under the top-level ``data/`` directory. This module is the one
place that resolves ``data/`` so that both the notebooks and the module ``main()``
entry points (``python -m mnn_torch.training --full``) agree on where results live,
regardless of the current working directory.

Because the package is installed editable, ``__file__`` resolves inside the real
repository ``src/`` tree, so the repository root -- and its sibling ``data/`` -- is
discoverable from the module location. Set ``MNN_DATA_DIR`` to override (e.g. for a
wheel install or CI where ``data/`` lives elsewhere).
"""
from __future__ import annotations

import os
from pathlib import Path

import numpy as np

__all__ = [
    "data_dir",
    "results_dir",
    "device_data_mat",
    "save_result",
    "load_result",
]


def data_dir() -> Path:
    """Absolute path to the package ``data/`` directory (created if absent).

    Honours the ``MNN_DATA_DIR`` environment variable; otherwise resolves to
    ``<repo-root>/data`` from this module's location (editable install).
    """
    env = os.environ.get("MNN_DATA_DIR")
    if env:
        d = Path(env).expanduser().resolve()
    else:
        # paths.py -> mnn_torch -> src -> <repo-root>
        d = Path(__file__).resolve().parents[2] / "data"
    d.mkdir(parents=True, exist_ok=True)
    return d


def results_dir() -> Path:
    """``data/results`` -- optional full-sweep caches used only when explicitly requested."""
    p = data_dir() / "results"
    p.mkdir(parents=True, exist_ok=True)
    return p


def device_data_mat() -> Path:
    """Path to the bundled SiOx multistate device fixture ``SiO_x-multistate-data.mat``.

    Resolves to ``data_dir()/SiO_x-multistate-data.mat`` when that exists. When
    ``MNN_DATA_DIR`` relocates the *writable* data dir (for MNIST downloads and
    result grids) the shipped fixture is not there, so we fall back to the
    repository-relative ``<repo-root>/data`` (discoverable from this module's
    location under the editable ``src/`` tree). The fixture is read-only, so it
    always ships with the repo even when results live elsewhere.
    """
    primary = data_dir() / "SiO_x-multistate-data.mat"
    if primary.exists():
        return primary
    repo_data = Path(__file__).resolve().parents[2] / "data" / "SiO_x-multistate-data.mat"
    if repo_data.exists():
        return repo_data
    return primary


def save_result(name: str, obj) -> Path:
    """Pickle-save ``obj`` to ``data/results/<name>`` (``.npy``), returning the path."""
    path = results_dir() / name
    np.save(path, obj, allow_pickle=True)
    return path


def load_result(name: str):
    """Load a result grid written by :func:`save_result` (``allow_pickle`` dict)."""
    return np.load(results_dir() / name, allow_pickle=True).item()
\end{lstlisting}

\begin{lstlisting}[caption={mnn\_torch/devices.py},label={lst:mnn-src-devices}]
import numpy as np
from scipy.io import loadmat


def load_SiOx_multistate(data_path) -> np.array:
    experimental_data = loadmat(data_path)["data"]
    experimental_data = np.flip(experimental_data, axis=2)
    experimental_data = np.transpose(experimental_data, (1, 2, 0))
    experimental_data = experimental_data[:2, :, :]

    return experimental_data


def load_SiOx_curves(
    experimental_data, max_voltage=0.5, voltage_step=0.005, clean_data=True
):
    num_points = int(max_voltage / voltage_step) + 1

    data = experimental_data[:, :, :num_points]
    if clean_data:
        data = clean_experimental_data(data)
    voltages = data[1, :, :]
    currents = data[0, :, :]

    return voltages, currents


def clean_experimental_data(experimental_data, threshold=0.1):
    accepted_curves = []
    for idx in range(experimental_data.shape[1]):
        curve = experimental_data[0, idx, :]
        d2y_dx2 = np.gradient(np.gradient(curve))
        ratio = d2y_dx2 / np.mean(curve)
        if ratio.max() < threshold:
            accepted_curves.append(idx)

    return experimental_data[:, accepted_curves, :]
\end{lstlisting}

\begin{lstlisting}[caption={mnn\_torch/effects.py},label={lst:mnn-src-effects}]
import numpy as np
from scipy.optimize import curve_fit
from scipy.stats import gaussian_kde, truncnorm
import scipy.constants as const
import torch

from mnn_torch.devices import load_SiOx_curves
from mnn_torch.utils import (
    sort_multiple_arrays,
    compute_multivariate_linear_regression,
    predict_with_multivariate_linear_regression,
)


def disturb_conductance_fixed(G, fixed_conductance, true_probability=0.5):
    """
    Disturbs the conductance by randomly replacing values with a fixed conductance.

    Args:
    - G: The original conductance array (tensor).
    - fixed_conductance: The fixed conductance value to replace randomly selected entries.
    - true_probability: The probability of replacing an entry with the fixed conductance.

    Returns:
    - G: The disturbed conductance array (tensor).
    - mask: The boolean mask indicating which values were replaced.
    """
    mask = torch.rand(*G.shape, device=G.device) < true_probability
    G = torch.where(mask, fixed_conductance, G)
    return G, mask


def disturb_conductance_device(G, G_on, G_off, true_probability=0.1):
    """
    Models stuck devices by probabilistically fixing their conductance to G_off or G_on.

    Args:
    - G (torch.Tensor): Original conductance tensor.
    - G_on (float or torch.Tensor): Maximum conductance value.
    - G_off (float or torch.Tensor): Minimum conductance value.
    - true_probability (float): Proportion of devices to be stuck (0 to 1).

    Returns:
    - torch.Tensor: Updated conductance tensor with probabilistic stuck behavior.
    """
    device = G.device
    G_on = torch.tensor(G_on, device=device, dtype=torch.float32)
    G_off = torch.tensor(G_off, device=device, dtype=torch.float32)

    # Calculate conductance range and median range
    conductance_range = G_on - G_off
    median_range = torch.median(conductance_range)

    # Identify potentially stuck devices based on range criteria
    stuck_mask = (G - G_off).abs() < (0.5 * median_range)

    # Generate a PDF for stuck values using KDE
    G_flat = G[~stuck_mask].detach().cpu().numpy()  # Non-stuck conductance values
    kde = gaussian_kde(G_flat)
    kde_samples = kde.resample(len(G_flat)).flatten()

    # Apply truncated normal distribution to mitigate bias near zero
    a, b = (0 - kde_samples.mean()) / kde_samples.std(), float("inf")
    truncated_samples = truncnorm(
        a, b, loc=kde_samples.mean(), scale=kde_samples.std()
    ).rvs(len(G_flat))
    truncated_samples = torch.tensor(
        truncated_samples, device=device, dtype=torch.float32
    )

    # Map stuck devices probabilistically to G_off or G_on
    stuck_values = torch.where(torch.rand_like(G, device=device) < 0.5, G_off, G_on)

    # Generate a random mask to select which devices are stuck
    random_stuck_mask = torch.rand_like(G, device=device) < true_probability

    # Incorporate device-to-device variability using a lognormal distribution
    # (Eq. lognormal). The lognormal model is defined on positive conductance, so
    # operate on |G| to stay valid when G carries a sign (e.g. conv kernels where
    # G = G_pos - G_neg); a positive floor avoids log(0)/NaN that would make the
    # sampling std non-finite. The original sign is restored after sampling.
    G_sign = torch.sign(G)
    G_mag = torch.clamp(torch.abs(G), min=1e-8)
    log_var = torch.clamp(torch.var(torch.log(G_mag)), min=0.0)
    log_std = torch.sqrt(log_var)
    log_mu = torch.log(G_mag) - (0.5 * log_var)
    # Sample per element: draw standard normals and scale by the (scalar) std so
    # the call is shape-safe across backends (MPS rejects a 0-dim std against a
    # full-shaped mean), then exponentiate. Equivalent to N(log_mu, log_std**2).
    noise = torch.randn_like(log_mu) * log_std
    variability_samples = G_sign * torch.exp(log_mu + noise)

    # Combine stuck and non-stuck conductance values out-of-place: stuck devices
    # take the fixed stuck value, the rest take the lognormal-perturbed value.
    # An out-of-place select (rather than in-place index assignment) is required
    # because G may be part of an autograd graph when the conductance feeds the
    # differentiable Poole-Frenkel forward map.
    G = torch.where(random_stuck_mask, stuck_values, variability_samples)

    return G


def sample_device_fault(shape, G_on, G_off, true_probability=0.1, sigma_g=0.5,
                        device="cpu", generator=None, stuck_polarity="split"):
    """Draw a *fixed* device-fault realisation, once, for a conductance array.

    This implements the manuscript fault protocol, in which the
    stuck-device pattern and the device-to-device (D2D) programming spread are
    sampled ONCE per training run and then held fixed, exactly as a physically
    programmed crossbar would behave: a device that is stuck, or programmed off
    its target, stays that way for the rest of training. This is distinct from
    :func:`disturb_conductance_device`, which re-rolls a fresh random fault on
    every forward pass and therefore injects a new, destabilising perturbation at
    every time step rather than a fixed hardware fault.

    The realisation is independent of the live weight values: it records *which*
    devices are stuck (``stuck_mask``), the conductance they are stuck at
    (``stuck_values``), and a per-device multiplicative lognormal D2D factor
    (``d2d_factor``, Eq. lognormal). It is applied later by
    :func:`apply_device_fault` to whatever conductance the current weights map to.

    The ``stuck_polarity`` argument follows the Joksas et al. nonideality
    taxonomy (``StuckAtGOff`` / ``StuckAtGOn``), which separates the two failure
    modes that the symmetric default blends together:

    * ``"off"`` -- stuck-at-G_off: dead / open synapses (cannot conduct). The
      presynaptic line goes silent; a firing-rate homeostat has nothing to act
      on, so this is the natural *negative control* for the homeostasis claim.
    * ``"on"``  -- stuck-at-G_on: shorted / saturated synapses that drive their
      postsynaptic line hyperactive. This is exactly the regime the sigmoidal
      firing-rate regulariser (Eq. homeo) is designed to suppress, so it is
      where homeostatic recovery should be largest.
    * ``"split"`` (default) -- each stuck device is G_off or G_on with equal
      probability, the original symmetric behaviour.

    Args:
    - shape (tuple): Shape of the conductance tensor the fault applies to.
    - G_on, G_off (float): Conductance bounds; stuck devices take one of these.
    - true_probability (float): Fraction of devices that are stuck.
    - sigma_g (float): Lognormal width of the D2D programming spread (Eq. lognormal).
    - device: Torch device for the realisation tensors.
    - generator (torch.Generator, optional): For reproducible realisations.
    - stuck_polarity (str): ``"split"`` | ``"off"`` | ``"on"`` (see above).

    Returns:
    - dict: {"stuck_mask", "stuck_values", "d2d_factor"} torch tensors of `shape`,
      all detached (a fixed hardware realisation, not a function of the weights).
    """
    G_on_t = torch.as_tensor(G_on, device=device, dtype=torch.float32)
    G_off_t = torch.as_tensor(G_off, device=device, dtype=torch.float32)

    def _rand(s):
        return torch.rand(s, device=device, generator=generator)

    def _randn(s):
        return torch.randn(s, device=device, generator=generator)

    stuck_mask = _rand(shape) < true_probability
    if stuck_polarity == "off":
        stuck_values = G_off_t.expand(shape).clone()
    elif stuck_polarity == "on":
        stuck_values = G_on_t.expand(shape).clone()
    elif stuck_polarity == "split":
        stuck_values = torch.where(_rand(shape) < 0.5, G_off_t, G_on_t)
    else:
        raise ValueError(
            f"stuck_polarity must be 'split', 'off' or 'on'; got {stuck_polarity!r}"
        )
    # Multiplicative lognormal D2D spread, mean-preserving (median at the target):
    # G_observed = G_target * exp(N(0, sigma_g^2)).
    d2d_factor = torch.exp(sigma_g * _randn(shape))
    return {
        "stuck_mask": stuck_mask.detach(),
        "stuck_values": stuck_values.detach(),
        "d2d_factor": d2d_factor.detach(),
    }


def apply_device_fault(G, fault):
    """Apply a *fixed* device-fault realisation (from :func:`sample_device_fault`).

    Stuck devices take their fixed stuck conductance; all other devices keep their
    programmed conductance scaled by the fixed per-device lognormal D2D factor. The
    operation is out-of-place so it is safe inside an autograd graph (the surviving
    devices' gradient still flows through ``G``), while the stuck devices and the
    D2D factors are constants that do not move during training.

    Args:
    - G (torch.Tensor): Live conductance tensor from the current weights.
    - fault (dict): Realisation from :func:`sample_device_fault`.

    Returns:
    - torch.Tensor: Conductance with the fixed fault realisation applied.
    """
    sign = torch.sign(G)
    perturbed = sign * torch.abs(G) * fault["d2d_factor"]
    return torch.where(fault["stuck_mask"], fault["stuck_values"], perturbed)


def compute_PooleFrenkel_current_torch(V, c, d_epsilon):
    r"""Differentiable Poole-Frenkel current (Torch), the forward map of the manuscript.

    Implements Eq. ``forward-g`` / Eq. ``pf``:

    .. math::
        I = c\, V \exp\!\left( \frac{2e}{k_B T} \sqrt{\frac{e |V|}{4\pi d\varepsilon}} \right)

    at ``T = 20 degC``, with the sign of ``V`` preserved so the map is valid for
    both bias polarities. All operations are Torch ops, so gradients flow through
    ``V``, ``c`` and ``d_epsilon``; this is the differentiable replacement for the
    NumPy :func:`compute_PooleFrenkel_current` used in offline fitting.

    Args:
        V (Tensor): Voltage.
        c (Tensor): Conductance-like prefactor (broadcastable to ``V``).
        d_epsilon (Tensor): Effective oxide-thickness--permittivity product
            (broadcastable to ``V``).

    Returns:
        Tensor: Poole-Frenkel current, same shape as the broadcast of the inputs.
    """
    e = const.elementary_charge
    kT = const.Boltzmann * (const.zero_Celsius + 20.0)
    V_abs = torch.abs(V)
    V_sign = torch.sign(V)
    # Manuscript Eq. pf: (2e/kT) * sqrt(e|V| / (4 pi d_epsilon)). This is
    # algebraically identical to the NumPy reference's (e/kT)*sqrt(e|V|/(pi d_eps))
    # since 2*sqrt(1/4) = 1. Clamp the radicand for numerical stability.
    term = torch.clamp(e * V_abs / (4.0 * np.pi * d_epsilon), min=1e-18)
    exponent = (2.0 * e / kT) * torch.sqrt(term)
    # Clamp the exponent: when a synapse's sampled conductance falls outside the
    # fitted resistance range, the regression can predict an unphysically small
    # d_epsilon and the raw exp() overflows to inf. Bounding the exponent keeps
    # the forward map finite and differentiable without distorting the in-range
    # response (exp(50) already exceeds any realistic current scale here).
    exponent = torch.clamp(exponent, max=50.0)
    return V_sign * c * V_abs * torch.exp(exponent)


def cholesky_factor_from_covariance(covariance_matrix, device="cpu", dtype=torch.float32):
    """Cholesky factor ``L`` of the PF residual covariance, with PD jitter.

    Precomputing this once (e.g. at layer construction) avoids recomputing a
    ``torch.linalg.cholesky`` on the constant covariance at every forward pass,
    which is the dominant redundant cost in the per-timestep PF sampling.
    """
    Sigma = torch.as_tensor(covariance_matrix, device=device, dtype=dtype)
    jitter = 1e-12 * torch.eye(2, device=device, dtype=dtype)
    return torch.linalg.cholesky(Sigma + jitter)


def sample_PooleFrenkel_parameters_torch(
    G, slopes, intercepts, covariance_matrix, stochastic=True, generator=None,
    cholesky_factor=None, residual=None, return_residual=False,
):
    r"""Sample per-synapse Poole-Frenkel parameters ``(c, d_epsilon)`` from conductance.

    Implements the bivariate regression of the manuscript (Eq. ``pf-regression``):

    .. math::
        [\ln c, \ln(d\varepsilon)]^\top = -\mathbf{m}\,\ln R + \mathbf{b} + \mathbf{E},
        \quad \mathbf{E} \sim \mathcal{N}_2(0, \Sigma)

    with ``R = 1/G``. The regression mean depends on the (trainable) conductance
    ``G`` and so carries gradient; the Gaussian residual ``E`` is treated as
    injected noise and is detached from the graph (stop-gradient), matching the
    manuscript's "Gradient flow" description. Returns ``c`` and ``d_epsilon`` as
    Torch tensors on ``G``'s device.

    Args:
        G (Tensor): Per-synapse conductance, any shape.
        slopes (array-like): Regression slopes ``m`` for ``[ln c, ln d_epsilon]``.
        intercepts (array-like): Regression intercepts ``b``.
        covariance_matrix (array-like): Residual covariance ``Sigma`` (2x2).
        stochastic (bool): If True, add the ``N(0, Sigma)`` residual; if False,
            return the regression mean only (deterministic).
        generator (torch.Generator, optional): RNG for reproducible sampling.
        cholesky_factor (Tensor, optional): Precomputed Cholesky factor of the
            residual covariance (see :func:`cholesky_factor_from_covariance`). When
            given, the per-call ``cholesky`` is skipped; ``covariance_matrix`` is
            then unused. Pass this from a layer that caches it once at construction.
        residual (Tensor, optional): A frozen ``(*G.shape, 2)`` device-to-device
            residual to reuse instead of drawing a fresh one. Models the variability
            as a fixed per-synapse property (sampled once per run), which also avoids
            the per-timestep ``randn``/matmul. ``cholesky_factor`` is then unused.
        return_residual (bool): If True, also return the residual that was used, so
            a caller can cache it for subsequent forward passes.

    Returns:
        tuple[Tensor, Tensor]: ``(c, d_epsilon)`` with the same shape as ``G``; or
        ``(c, d_epsilon, residual)`` when ``return_residual`` is True.
    """
    device = G.device
    dtype = G.dtype
    m = torch.as_tensor(slopes, device=device, dtype=dtype)
    b = torch.as_tensor(intercepts, device=device, dtype=dtype)

    ln_R = torch.log(1.0 / torch.clamp(G, min=1e-12))
    # mean of [ln c, ln d_epsilon]: shape (*G.shape, 2).
    # NOTE: slopes/intercepts come from scipy.linregress (compute_multivariate_
    # linear_regression), whose convention is y = slope*x + intercept, so the
    # prediction uses +slopes here. This is the signed form of the manuscript's
    # Eq. pf-regression (-m ln R + b), where the code's slope already carries the
    # sign; the +convention yields physical d_epsilon ~1e-18 matching the fit.
    mean = ln_R.unsqueeze(-1) * m + b

    drawn_residual = None
    if stochastic:
        if residual is not None:
            # Frozen device-to-device residual (sampled once per run): the
            # variability is a fixed property of each fabricated synapse, so it is
            # not re-rolled every timestep. The weight-dependent mean still updates,
            # so gradients and training-time conductance changes are preserved.
            drawn_residual = residual.to(device=device, dtype=dtype)
        else:
            # Use a precomputed Cholesky factor when supplied (constant across
            # forward passes); otherwise compute it from the covariance. Computing
            # it once at layer construction removes a linalg op from the timestep loop.
            if cholesky_factor is not None:
                L = cholesky_factor.to(device=device, dtype=dtype)
            else:
                L = cholesky_factor_from_covariance(covariance_matrix, device, dtype)
            z = torch.randn(*G.shape, 2, device=device, dtype=dtype, generator=generator)
            drawn_residual = z @ L.T
        ln_params = mean + drawn_residual.detach()  # stop-gradient on the noise
    else:
        ln_params = mean

    c = torch.exp(ln_params[..., 0])
    d_epsilon = torch.exp(ln_params[..., 1])
    if return_residual:
        return c, d_epsilon, drawn_residual
    return c, d_epsilon


def compute_PooleFrenkel_current(V, c, d_epsilon):
    """
    Computes the Poole-Frenkel current for a given voltage.

    Args:
    - V: The voltage array.
    - c: A coefficient related to the current.
    - d_epsilon: A parameter related to the material's dielectric properties.

    Returns:
    - I: The calculated current array.
    """
    V_abs = np.abs(V)
    V_sign = np.sign(V)
    term = np.maximum(const.elementary_charge * V_abs / (const.pi * d_epsilon), 1e-18)
    I = (
        V_sign
        * c
        * V_abs
        * np.exp(
            const.elementary_charge
            * np.sqrt(term)
            / (const.Boltzmann * (const.zero_Celsius + 20.0))
        )
    )
    return I


def compute_PooleFrenkel_total_current(V, G, slopes, intercepts, covariance_matrix):
    """
    Computes the total Poole-Frenkel current, including deviations based on the regression parameters.

    Args:
    - V: The voltage array.
    - G: The conductance array.
    - slopes: The slopes from the regression model.
    - intercepts: The intercepts from the regression model.
    - covariance_matrix: The covariance matrix for the regression.

    Returns:
    - I: The total current.
    - I_ind: The individual current components.
    """
    R = 1 / G
    ln_R = np.log(R)

    # Predict regression results for ln(R)
    fit_data = predict_with_multivariate_linear_regression(
        ln_R, slopes, intercepts, covariance_matrix
    )

    # Extract coefficients for the Poole-Frenkel equation
    c = np.exp(fit_data[0])
    d_epsilon = np.exp(fit_data[1])

    # Compute individual currents
    I_ind = compute_PooleFrenkel_current(V, c, d_epsilon)

    # Sum currents along the bit lines
    I = np.sum(I_ind, axis=1)
    return I, I_ind


def compute_PooleFrenkel_relationship(V, I, voltage_step=0.005, ref_voltage=0.1):
    """
    Computes the relationship between voltage and current for the Poole-Frenkel model.

    Args:
    - V: Voltage array.
    - I: Current array.
    - voltage_step: Step size for voltage.
    - ref_voltage: The reference voltage to compute resistance.

    Returns:
    - R, c, d_epsilon: The computed resistance, c, and d_epsilon values.
    - V, I: The sorted voltage and current arrays.
    """
    num_curves = V.shape[0]
    R = np.zeros(num_curves)
    c = np.zeros(num_curves)
    d_epsilon = np.zeros(num_curves)

    ref_idx = int(ref_voltage / voltage_step)

    for idx in range(num_curves):
        v = V[idx, :]
        i = I[idx, :]

        # Compute resistance at the reference voltage
        if i[ref_idx] == 0:  # Avoid division by zero
            R[idx] = np.inf
        else:
            R[idx] = v[ref_idx] / i[ref_idx]

        # Fit the Poole-Frenkel current model
        try:
            popt, _ = curve_fit(
                compute_PooleFrenkel_current,
                v,
                i,
                p0=[1e-5, 1e-16],
                bounds=([1e-8, 1e-18], [1e-3, 1e-14]),
                maxfev=10000,  # Increase iterations if needed
            )
            c[idx] = popt[0]
            d_epsilon[idx] = popt[1]
        except RuntimeError as e:
            print(f"Curve fit failed for curve {idx}: {e}")
            c[idx], d_epsilon[idx] = np.nan, np.nan

    # Sort the results
    R, c, d_epsilon, V, I = sort_multiple_arrays(R, c, d_epsilon, V, I)
    return R, c, d_epsilon, V, I


def compute_PooleFrenkel_parameters(
    experimental_data, high_resistance_state=False, ratio=5
):
    """
    Computes the Poole-Frenkel model parameters based on experimental data.

    Args:
    - experimental_data: The data used for the calculations.
    - high_resistance_state: Boolean indicating whether to use a high resistance state.
    - ratio: The ratio of the conductance in "on" and "off" states.

    Returns:
    - G_off, G_on: The conductance in the "off" and "on" states.
    - R, c, d_epsilon: The computed resistance, c, and d_epsilon values.
    """
    V, I = load_SiOx_curves(experimental_data)
    R, c, d_epsilon, _, _ = compute_PooleFrenkel_relationship(V, I)

    # Convert results to float32 for consistency
    R = R.astype(np.float32)
    c = c.astype(np.float32)
    d_epsilon = d_epsilon.astype(np.float32)

    # Compute G_on and G_off based on the resistance state
    if high_resistance_state:
        G_off = 1 / R[-1]
        G_on = G_off * ratio
    else:
        G_on = 1 / R[0]
        G_off = G_on / ratio

    return G_off, G_on, R, c, d_epsilon


def compute_PooleFrenkel_regression_parameters(
    R, c, d_epsilon, high_resistance_state=False
):
    """
    Computes the regression parameters for the Poole-Frenkel model.

    Args:
    - R: The resistance array.
    - c: The coefficient array.
    - d_epsilon: The dielectric constant array.
    - high_resistance_state: Boolean indicating the state.

    Returns:
    - slopes, intercepts, covariance_matrix: The regression parameters.
    """
    sep_idx = np.searchsorted(
        R, const.physical_constants["inverse of conductance quantum"][0]
    )

    if high_resistance_state:
        x = np.log(R[sep_idx:])
        y_1 = np.log(c[sep_idx:])
        y_2 = np.log(d_epsilon[sep_idx:])
    else:
        x = np.log(R[:sep_idx])
        y_1 = np.log(c[:sep_idx])
        y_2 = np.log(d_epsilon[:sep_idx])

    # Compute regression parameters
    slopes, intercepts, covariance_matrix = compute_multivariate_linear_regression(
        x, y_1, y_2
    )

    return slopes, intercepts, covariance_matrix
\end{lstlisting}

\begin{lstlisting}[caption={mnn\_torch/layers.py},label={lst:mnn-src-layers}]
import torch
import torch.nn as nn
import numpy as np

from mnn_torch.effects import (
    disturb_conductance_fixed,
    disturb_conductance_device,
    sample_device_fault,
    apply_device_fault,
    compute_PooleFrenkel_current_torch,
    sample_PooleFrenkel_parameters_torch,
    cholesky_factor_from_covariance,
)


class MemristiveLinearLayer(nn.Module):
    """
    Custom memristive linear layer.
    This layer simulates the behavior of memristors, optionally applying
    disturbances and non-ideal effects based on the provided configuration.
    """

    def __init__(self, size_in, size_out, memristive_config):
        """
        Initializes the memristive linear layer.

        Args:
        - size_in (int): Number of input features.
        - size_out (int): Number of output features.
        - memristive_config (dict): Configuration dictionary containing:
            - 'ideal': If True, the layer behaves like a standard linear layer.
            - 'disturb_conductance': Whether to disturb conductance values.
            - 'disturb_mode': The mode of disturbance ('fixed' or 'device').
            - 'disturbance_probability': The probability for disturbance.
            - 'G_off', 'G_on': Minimum and maximum conductance values.
            - 'R', 'c', 'd_epsilon', 'k_V': Parameters for non-ideal effects.
        """
        super().__init__()
        self.size_in = size_in
        self.size_out = size_out

        # Load configuration
        self.disturb_conductance = memristive_config.get("disturb_conductance", False)
        self.disturb_mode = memristive_config.get("disturb_mode", "fixed")
        self.disturbance_probability = memristive_config.get(
            "disturbance_probability", 0.1
        )
        # Lognormal D2D width for the fixed-fault ('device_fixed') protocol.
        self.disturbance_sigma_g = memristive_config.get("disturbance_sigma_g", 0.5)
        # Stuck-fault polarity for 'device_fixed': 'split' (G_off or G_on, symmetric
        # default), 'off' (dead/open synapses) or 'on' (shorted/saturated synapses).
        # Load-bearing for the homeostasis stuck-on vs stuck-off mechanism test.
        self.disturbance_stuck_polarity = memristive_config.get(
            "disturbance_stuck_polarity", "split"
        )
        # Cache for the fixed device-fault realisation (Section 6.3 protocol):
        # sampled once on first forward and reused for the rest of the run.
        self._device_fault = None
        self.homeostasis_dropout = memristive_config.get("homeostasis_dropout", False)

        # Initialize weights and biases
        self.weights = nn.Parameter(torch.Tensor(size_out, size_in))
        self.bias = nn.Parameter(torch.Tensor(size_out))
        stdv = 1 / np.sqrt(self.size_in)
        nn.init.normal_(self.weights, mean=0, std=stdv)
        nn.init.constant_(self.bias, 0)

        self.G_off = memristive_config["G_off"]
        self.G_on = memristive_config["G_on"]
        self.R = memristive_config["R"]
        self.c = memristive_config["c"]
        self.d_epsilon = memristive_config["d_epsilon"]
        self.k_V = memristive_config["k_V"]

        # Poole-Frenkel nonlinear forward map (manuscript Eq. forward-g).
        # When enabled, the per-synapse current is computed from the differentiable
        # PF model with (c, d_epsilon) sampled from the bivariate regression of
        # Eq. pf-regression, instead of the linear Ohmic product I = V * G.
        self.nonlinear = memristive_config.get("nonlinear", False)
        self.pf_slopes = memristive_config.get("pf_slopes")
        self.pf_intercepts = memristive_config.get("pf_intercepts")
        self.pf_covariance = memristive_config.get("pf_covariance")
        self.pf_stochastic = memristive_config.get("pf_stochastic", True)
        if self.nonlinear and (
            self.pf_slopes is None
            or self.pf_intercepts is None
            or self.pf_covariance is None
        ):
            raise ValueError(
                "nonlinear=True requires 'pf_slopes', 'pf_intercepts' and "
                "'pf_covariance' in memristive_config (see "
                "compute_PooleFrenkel_regression_parameters)."
            )
        # Precompute the Cholesky factor of the PF residual covariance once, so the
        # per-timestep PF sampling skips a redundant linalg.cholesky on a constant.
        self._pf_chol = (
            cholesky_factor_from_covariance(self.pf_covariance)
            if (self.nonlinear and self.pf_stochastic)
            else None
        )
        # When True, the PF device-to-device residual is sampled ONCE and frozen
        # for the run (a fixed per-synapse property), rather than re-rolled every
        # timestep. This is both faster (no per-step randn/matmul) and the more
        # physically faithful reading of device-to-device variability.
        self.pf_frozen_variability = memristive_config.get(
            "pf_frozen_variability", False
        )
        self._pf_residual = None

    def forward(self, x):
        """
        Forward pass for the memristive linear layer.

        Args:
        - x (Tensor): Input tensor of shape (batch_size, size_in).

        Returns:
        - Tensor: Output tensor of shape (batch_size, size_out).
        """
        device = x.device

        # Include bias in input
        inputs = torch.cat([x, torch.ones([x.shape[0], 1], device=device)], dim=1)
        bias = torch.unsqueeze(self.bias, dim=0).to(device)
        weights_and_bias = torch.cat([self.weights.t(), bias], dim=0)

        # Compute the memristive outputs
        outputs = self.memristive_outputs(inputs, weights_and_bias)
        return outputs

    def memristive_outputs(self, x, weights_and_bias):
        """
        Computes the memristive layer outputs based on voltage, conductance,
        and current relationships.

        Args:
        - x (Tensor): Input tensor with bias included, shape (batch_size, size_in + 1).
        - weights_and_bias (Tensor): Weights and bias tensor, shape (size_in + 1, size_out).

        Returns:
        - Tensor: Output tensor of shape (batch_size, size_out).
        """
        # Voltage calculation
        V = self.k_V * x

        # Conductance scaling factors
        max_wb = torch.max(torch.abs(weights_and_bias))
        k_G = (self.G_on - self.G_off) / max_wb
        k_I = self.k_V * k_G
        G_eff = k_G * weights_and_bias

        # Map weights to positive and negative conductances
        G_pos = self.G_off + torch.maximum(G_eff, torch.tensor(0.0, device=x.device))
        G_neg = self.G_off - torch.minimum(G_eff, torch.tensor(0.0, device=x.device))
        G = torch.cat((G_pos[:, :, None], G_neg[:, :, None]), dim=-1).reshape(
            G_pos.size(0), -1
        )

        # Disturb conductance based on the selected mode (stuck-at / device faults).
        # Homeostasis is applied at the spike level by the model's regulariser,
        # not here at the conductance level, so the two concerns stay separate.
        if self.disturb_conductance:
            if self.disturb_mode == "fixed":
                G, _mask = disturb_conductance_fixed(
                    G, self.G_on, true_probability=self.disturbance_probability
                )
            elif self.disturb_mode == "device":
                G = disturb_conductance_device(
                    G,
                    self.G_on,
                    self.G_off,
                    true_probability=self.disturbance_probability,
                )
            elif self.disturb_mode == "device_fixed":
                # Section 6.3 protocol: sample the stuck mask + D2D spread ONCE
                # (on first forward, when G's shape is known) and hold it fixed.
                if self._device_fault is None:
                    self._device_fault = sample_device_fault(
                        G.shape, self.G_on, self.G_off,
                        true_probability=self.disturbance_probability,
                        sigma_g=self.disturbance_sigma_g, device=G.device,
                        stuck_polarity=self.disturbance_stuck_polarity,
                    )
                G = apply_device_fault(G, self._device_fault)

        # Compute per-synapse current.
        # V: (batch, size_in+1); G: (size_in+1, 2*size_out) interleaved pos/neg.
        V_b = torch.unsqueeze(V, dim=-1)  # (batch, size_in+1, 1)
        G_b = torch.unsqueeze(G, dim=0)  # (1, size_in+1, 2*size_out)
        if self.nonlinear:
            # Poole-Frenkel nonlinear map (Eq. forward-g): sample (c, d_epsilon)
            # per synapse from its conductance, then I = PF(V, c, d_epsilon).
            if self.pf_frozen_variability and self.pf_stochastic:
                # Sample the device-to-device residual once (on first forward, when
                # G's broadcast shape is known) and reuse it thereafter.
                c, d_epsilon, res = sample_PooleFrenkel_parameters_torch(
                    G_b,
                    self.pf_slopes,
                    self.pf_intercepts,
                    self.pf_covariance,
                    stochastic=True,
                    cholesky_factor=self._pf_chol,
                    residual=self._pf_residual,
                    return_residual=True,
                )
                if self._pf_residual is None:
                    self._pf_residual = res.detach()
            else:
                c, d_epsilon = sample_PooleFrenkel_parameters_torch(
                    G_b,
                    self.pf_slopes,
                    self.pf_intercepts,
                    self.pf_covariance,
                    stochastic=self.pf_stochastic,
                    cholesky_factor=self._pf_chol,
                )
            I_ind = compute_PooleFrenkel_current_torch(V_b, c, d_epsilon)
        else:
            # Linear Ohmic map: I = V * G.
            I_ind = V_b * G_b
        I = torch.sum(I_ind, dim=1)
        I_total = I[:, 0::2] - I[:, 1::2]

        # Output calculation
        y = I_total / k_I
        return y


class MemristiveConv2d(nn.Module):
    """
    Custom memristive 2D convolutional layer.
    This layer simulates the behavior of memristors, optionally applying
    disturbances and non-ideal effects based on the provided configuration.
    """

    def __init__(
        self,
        in_channels,
        out_channels,
        kernel_size,
        stride=1,
        padding=0,
        memristive_config=None,
    ):
        """
        Initializes the memristive convolutional layer.
        """
        super().__init__()
        self.in_channels = in_channels
        self.out_channels = out_channels
        self.kernel_size = self._to_tuple(kernel_size)
        self.stride = self._to_tuple(stride)
        self.padding = self._to_tuple(padding)

        # Load configuration
        self.ideal = memristive_config.get("ideal", True)
        self.disturb_conductance = memristive_config.get("disturb_conductance", False)
        self.disturb_mode = memristive_config.get("disturb_mode", "fixed")
        self.disturbance_probability = memristive_config.get(
            "disturbance_probability", 0.1
        )
        self.disturbance_sigma_g = memristive_config.get("disturbance_sigma_g", 0.5)
        # Stuck-fault polarity for 'device_fixed' (see effects.sample_device_fault):
        # 'split' (symmetric), 'off' (dead synapses) or 'on' (saturated synapses).
        self.disturbance_stuck_polarity = memristive_config.get(
            "disturbance_stuck_polarity", "split"
        )
        self._device_fault = None
        self.homeostasis_dropout = memristive_config.get("homeostasis_dropout", False)

        # Initialize weights and biases
        self.weight = nn.Parameter(
            torch.Tensor(out_channels, in_channels, *self.kernel_size)
        )
        self.bias = nn.Parameter(torch.Tensor(out_channels))
        stdv = 1 / np.sqrt(in_channels * np.prod(self.kernel_size))
        nn.init.normal_(self.weight, mean=0, std=stdv)
        nn.init.constant_(self.bias, 0)

        # Load non-ideal parameters from config
        self.G_off = memristive_config["G_off"]
        self.G_on = memristive_config["G_on"]
        self.R = memristive_config["R"]
        self.c = memristive_config["c"]
        self.d_epsilon = memristive_config["d_epsilon"]
        self.k_V = memristive_config["k_V"]

        # Poole-Frenkel nonlinear forward map for the convolution (beyond Joksas et al.,
        # who applied PF to dense layers only). A conv KERNEL weight is one physical
        # device reused across every spatial position, so the stuck-fault mask and the
        # PF device-to-device residual are sampled ONCE over the (K, 2*out_channels)
        # conductance grid and broadcast over batch and space.
        self.nonlinear = memristive_config.get("nonlinear", False)
        self.pf_slopes = memristive_config.get("pf_slopes")
        self.pf_intercepts = memristive_config.get("pf_intercepts")
        self.pf_covariance = memristive_config.get("pf_covariance")
        self.pf_stochastic = memristive_config.get("pf_stochastic", True)
        if self.nonlinear and (
            self.pf_slopes is None
            or self.pf_intercepts is None
            or self.pf_covariance is None
        ):
            raise ValueError(
                "nonlinear=True requires 'pf_slopes', 'pf_intercepts' and "
                "'pf_covariance' in memristive_config (see "
                "compute_PooleFrenkel_regression_parameters)."
            )
        self._pf_chol = (
            cholesky_factor_from_covariance(self.pf_covariance)
            if (self.nonlinear and self.pf_stochastic)
            else None
        )
        self.pf_frozen_variability = memristive_config.get(
            "pf_frozen_variability", False
        )
        self._pf_residual = None
        # Memory bounding for the (B, L, K, 2*out) PF current tensor (the slow im2col
        # path): chunk over spatial L; pf_checkpoint bounds the backward peak.
        self.pf_chunk_L = memristive_config.get("pf_chunk_L", None)
        self.pf_checkpoint = memristive_config.get("pf_checkpoint", False)

    def _to_tuple(self, value):
        """Ensures kernel size and padding are tuples."""
        return value if isinstance(value, tuple) else (value, value)

    def forward(self, x):
        """
        Forward pass for the memristive convolutional layer.

        Args:
        - x (Tensor): Input tensor of shape (batch_size, in_channels, height, width).

        Returns:
        - Tensor: Output tensor (batch_size, out_channels, new_height, new_width).
        """
        if self.ideal:
            return nn.functional.conv2d(
                x, self.weight, self.bias, self.stride, self.padding
            )
        # Both the linear-memristive and PF-nonlinear paths go through the same im2col
        # current-accumulation pipeline; only the per-element current model differs.
        return self._memristive_conv(x)

    def _conductance_grid(self, device):
        """Map the kernel to the interleaved pos/neg conductance grid and apply the
        sampled-once stuck-fault realisation. Returns ``(G, k_I)`` with ``G`` of shape
        ``(K, 2*out_channels)`` (K = in_channels*kh*kw), mirroring the dense layer.
        Gradient flows through ``self.weight``."""
        Cout = self.out_channels
        K = self.in_channels * self.kernel_size[0] * self.kernel_size[1]
        Wmat = self.weight.reshape(Cout, K).t()            # (K, Cout)
        # max over the KERNEL only (bias added additively in output units), so a conv
        # weight and a dense weight of equal value see identical device physics and the
        # 1x1-conv == dense identity is exact.
        max_w = torch.max(torch.abs(self.weight))
        k_G = (self.G_on - self.G_off) / max_w
        k_I = self.k_V * k_G
        G_eff = k_G * Wmat
        zero = torch.tensor(0.0, device=device)
        G_pos = self.G_off + torch.maximum(G_eff, zero)
        G_neg = self.G_off - torch.minimum(G_eff, zero)
        G = torch.cat((G_pos[:, :, None], G_neg[:, :, None]), dim=-1).reshape(
            K, 2 * Cout
        )

        if self.disturb_conductance:
            if self.disturb_mode == "fixed":
                G, _mask = disturb_conductance_fixed(
                    G, self.G_on, true_probability=self.disturbance_probability
                )
            elif self.disturb_mode == "device":
                G = disturb_conductance_device(
                    G, self.G_on, self.G_off,
                    true_probability=self.disturbance_probability,
                )
            elif self.disturb_mode == "device_fixed":
                if self._device_fault is None:
                    self._device_fault = sample_device_fault(
                        G.shape, self.G_on, self.G_off,
                        true_probability=self.disturbance_probability,
                        sigma_g=self.disturbance_sigma_g, device=G.device,
                        stuck_polarity=self.disturbance_stuck_polarity,
                    )
                G = apply_device_fault(G, self._device_fault)
        return G, k_I

    def _pf_params(self, G):
        """Per-device PF parameters ``(c, d_epsilon)`` sampled once on the static
        ``(1, 1, K, 2*Cout)`` grid; frozen residual cached so later forwards reproduce
        it (the mean still carries gradient through ``G``)."""
        G_b = G[None, None, :, :]
        if self.pf_frozen_variability and self.pf_stochastic:
            c, d_epsilon, res = sample_PooleFrenkel_parameters_torch(
                G_b, self.pf_slopes, self.pf_intercepts, self.pf_covariance,
                stochastic=True, cholesky_factor=self._pf_chol,
                residual=self._pf_residual, return_residual=True,
            )
            if self._pf_residual is None:
                self._pf_residual = res.detach()
        else:
            c, d_epsilon = sample_PooleFrenkel_parameters_torch(
                G_b, self.pf_slopes, self.pf_intercepts, self.pf_covariance,
                stochastic=self.pf_stochastic, cholesky_factor=self._pf_chol,
            )
        return c, d_epsilon

    def _current_tile(self, V_tile, G, c, d_epsilon, k_I):
        """Per-spatial-tile current over (B, t, K, 2*Cout), summed over K, pos-neg, /k_I."""
        if self.nonlinear:
            I_ind = compute_PooleFrenkel_current_torch(V_tile, c, d_epsilon)
        else:
            I_ind = V_tile * G[None, None, :, :]
        I = torch.sum(I_ind, dim=2)
        return (I[..., 0::2] - I[..., 1::2]) / k_I

    def _memristive_conv(self, x):
        B = x.shape[0]
        Cout = self.out_channels
        kh, kw = self.kernel_size
        sh, sw = self.stride
        ph, pw = self.padding
        oh = (x.shape[2] + 2 * ph - kh) // sh + 1
        ow = (x.shape[3] + 2 * pw - kw) // sw + 1

        G, k_I = self._conductance_grid(x.device)          # (K, 2*Cout)
        c = d_epsilon = None
        if self.nonlinear:
            c, d_epsilon = self._pf_params(G)

            # Binary-spike fast path. When the input is a spike train (every element 0
            # or 1 -- MCSNN conv2 fed by a LIF spike output), PF(k_V*0)=0 and
            # PF(k_V*1)=I_on, so the per-synapse current folds into a fixed effective
            # kernel and the whole conv becomes a single F.conv2d: the exp/sqrt runs ONCE
            # over (K,2*Cout) instead of over (B,L,K,2*Cout). Measured ~19x faster than
            # the im2col path at conv2 scale. Gradient flows through I_on->(c,d_eps)->G
            # ->weight. The runtime guard MUST fall through to full PF whenever the
            # homeostasis regulariser has gated the input to continuous (0,1] values --
            # those are exactly the timesteps the homeostasis study measures.
            if torch.all((x == 0) | (x == 1)):
                K = G.shape[0]
                k_V_full = torch.full((1, 1, K, 2 * Cout), self.k_V,
                                      device=x.device, dtype=x.dtype)
                I_on = compute_PooleFrenkel_current_torch(
                    k_V_full, c, d_epsilon).reshape(K, 2 * Cout)
                W_eff = (I_on[:, 0::2] - I_on[:, 1::2]) / k_I       # (K, Cout)
                W_eff_kernel = W_eff.t().reshape(Cout, self.in_channels, kh, kw)
                out = nn.functional.conv2d(
                    x, W_eff_kernel, None, self.stride, self.padding)
                return out + self.bias.view(1, Cout, 1, 1)

        # im2col full-PF path: input patches (B, K, L), L = oh*ow output positions.
        patches = nn.functional.unfold(
            x, kernel_size=self.kernel_size, stride=self.stride, padding=self.padding
        )
        V = self.k_V * patches                             # (B, K, L)
        L = V.shape[-1]

        Lc = self.pf_chunk_L or L
        outs = []
        for s in range(0, L, Lc):
            V_tile = V[:, :, s:s + Lc].permute(0, 2, 1).unsqueeze(-1)  # (B, t, K, 1)
            if self.pf_checkpoint and self.training and self.nonlinear:
                tile = torch.utils.checkpoint.checkpoint(
                    self._current_tile, V_tile, G, c, d_epsilon, k_I,
                    use_reentrant=False,
                )
            else:
                tile = self._current_tile(V_tile, G, c, d_epsilon, k_I)
            outs.append(tile)
        y_cols = torch.cat(outs, dim=1)                    # (B, L, Cout)

        out = y_cols.permute(0, 2, 1).reshape(B, Cout, oh, ow)
        return out + self.bias.view(1, Cout, 1, 1)


class HomeostaticRegulariser(nn.Module):
    r"""Firing-rate-dependent homeostatic regulariser (manuscript Eq. homeo).

    Modulates the effective conductance of a presynaptic line as a smooth,
    differentiable sigmoidal function of its firing rate ``r_in`` measured over a
    sliding window:

    .. math::
        \frac{G_\mathrm{eff}}{G_\mathrm{max}} = \frac{1}{1 + (r_\mathrm{in}/r_\mathrm{th})^{n}}

    applied as a multiplicative gate on the presynaptic spike train. The gate is
    ~1 for ``r_in << r_th`` (normal regime untouched) and decays as
    ``r_in ** -n`` for ``r_in >> r_th`` (hyperactive lines suppressed). It is
    row-local (depends only on the presynaptic line, no postsynaptic feedback) and
    differentiable in ``r_in``, ``r_th`` and ``n``, so it composes with
    surrogate-gradient BPTT and admits gradient-based hyperparameter tuning.

    This replaces the binary, non-differentiable :class:`HomeostasisDropout`,
    which is retained only for backward compatibility.

    Args:
        r_th (float): Firing-rate threshold at which depression begins, expressed
            as a fraction of time steps spiked in the window (``[0, 1]``).
        n (float): Sharpness exponent of the sigmoidal transition. Defaults to the
            parameter-free derived value ``n=1`` (single screening state, the steady
            state of the Ch5 fill-emission trap balance, Eq. 5.15); ``n>1`` is a
            cascade-sharpened probe, not the grounded default.
    """

    def __init__(self, r_th=0.8, n=1.0):
        super().__init__()
        self.r_th = r_th
        self.n = n

    def gate(self, spk_window):
        """Per-neuron conductance gate in ``(0, 1]`` from a window of spikes.

        Args:
        - spk_window (Tensor): Recent presynaptic spikes, shape
          ``(window, batch, *features)``.

        Returns:
        - Tensor: Gate of shape ``(batch, *features)``.
        """
        r_in = spk_window.mean(dim=0)
        ratio = torch.clamp(r_in / self.r_th, min=0.0)
        return 1.0 / (1.0 + ratio.pow(self.n))

    def forward(self, spk_window):
        """Scale the current-step spikes by the windowed homeostatic gate.

        Args:
        - spk_window (Tensor): Recent presynaptic spikes, shape
          ``(window, batch, *features)``; the last entry is the current step.

        Returns:
        - Tensor: Gated current-step spikes, shape ``(batch, *features)``.
        """
        return spk_window[-1] * self.gate(spk_window)


class HomeostasisDropout(nn.Module):
    def __init__(self):
        """
        A custom layer that drops the spike output if it has been spiking continuously
        over all time steps (across the entire sequence).

        Legacy binary mechanism, retained for backward compatibility. The
        firing-rate-dependent :class:`HomeostaticRegulariser` implements the
        manuscript's Eq. homeo and is the default in the models.
        """
        super(HomeostasisDropout, self).__init__()

    def forward(self, spk_rec):
        """
        Forward pass that checks if the spikes are repeated across all time steps
        and sets those to zero if they are continuously '1' over the entire sequence.

        Args:
        - spk_rec: A tensor containing the spikes over time (num_steps, batch_size, num_features).

        Returns:
        - A tensor with the spikes set to zero if they were '1' continuously across the entire sequence.
        """

        # The shape of spk_rec is (num_steps, batch_size, num_features)
        num_steps, batch_size, *feature_dims = spk_rec.shape

        # Check for continuous spikes across all time steps
        # Compare each time step with the next one to find continuous sequences of 1s
        continuous_spikes = (spk_rec[1:] == 1) & (
            spk_rec[:-1] == 1
        )  # Shape: (num_steps-1, batch_size, num_features)

        # We want to drop the spikes that were continuous across all time steps
        # The spike is dropped if it was '1' in all steps
        drop_mask = torch.cat(
            [
                torch.zeros(1, batch_size, *feature_dims, device=spk_rec.device),
                continuous_spikes,
            ],
            dim=0,
        )

        # Drop the spikes that were continuous across all time steps by applying the drop mask
        spk_rec = spk_rec * (1 - drop_mask).float()

        # Return only the latest step after applying the drop mask
        return spk_rec[-1]  # Select the last step in the sequence
\end{lstlisting}

\begin{lstlisting}[caption={mnn\_torch/models.py},label={lst:mnn-src-models}]
import torch
import torch.nn as nn
import snntorch as snn
import torch.nn.functional as F
from mnn_torch.layers import (
    MemristiveLinearLayer,
    MemristiveConv2d,
    HomeostasisDropout,
    HomeostaticRegulariser,
)


def beta_from_tau_leak(tau_leak, dt):
    r"""Map a measured device retention ``tau_leak`` to the LIF membrane decay ``beta``.

    A leaky integrator with time constant :math:`\tau` decays over one timestep of
    length :math:`dt` by the factor :math:`\beta = \exp(-dt/\tau)`. This is the SAME
    exponentially-leaky state, and the SAME :math:`\tau_{\mathrm{leak}}` (trap-discharge
    time, :math:`R_{\mathrm{leak}}C`), used as the eligibility
    trace in the reward-modulated learning study. Using it here makes the neuron's temporal
    membrane the device's own relaxation: one measured material timescale instantiates
    both the reinforcement-learning eligibility trace and the spiking-classifier's temporal integrator --
    the leaky-integrator primitive the two share. ``beta`` is clamped to ``[0, 1)`` as
    snnTorch requires.
    """
    import math
    beta = math.exp(-float(dt) / float(tau_leak))
    return min(max(beta, 0.0), 1.0 - 1e-6)


class DeviceLeaky(snn.Leaky):
    """A LIF neuron whose membrane leak is set by the MEASURED device retention
    ``tau_leak`` rather than a free hyperparameter: ``beta = exp(-dt/tau_leak)``
    (see :func:`beta_from_tau_leak`). Semantically identical to ``snn.Leaky`` -- it
    only fixes where ``beta`` comes from -- so it is a drop-in whose membrane dynamics
    are device-grounded, matching the eligibility-trace substrate used in the
    companion reinforcement-learning experiments."""

    def __init__(self, tau_leak, dt, spike_grad=None, **kwargs):
        super().__init__(beta=beta_from_tau_leak(tau_leak, dt),
                         spike_grad=spike_grad, **kwargs)
        self.tau_leak = float(tau_leak)
        self.dt = float(dt)


class AdaptiveLeaky(nn.Module):
    r"""Adaptive leaky integrate-and-fire (ALIF) neuron -- the LSNN working-memory cell
    of Bellec et al. (NeurIPS 2018 / Nat. Commun. 2020).

    A plain LIF neuron has one state, the membrane, and a FIXED threshold. The ALIF
    neuron adds a SECOND, slow per-neuron state: an adaptive threshold. Each spike
    raises the threshold; it then decays with its own long time constant
    :math:`\tau_a \gg \tau_m`:

    .. math::
        v_{t}   &= \beta\, v_{t-1} + I_t - z_{t-1}\,\vartheta_{t-1} \\
        b_{t}   &= \rho\, b_{t-1} + z_{t-1},\quad \rho = e^{-\Delta t/\tau_a} \\
        \vartheta_t &= v_{\mathrm{th}} + \gamma\, b_t \\
        z_t     &= H(v_t - \vartheta_t)

    Information written into the slow adaptation state :math:`b` persists for
    :math:`\sim\tau_a`, which is what gives the LSNN its ``long short-term memory''.
    This is the ENGINEERED-variable route to a slow timescale, and the baseline against
    which :class:`DeviceLeaky` is compared: the device supplies the same slow timescale
    from the material leak alone (``beta = exp(-dt/tau_leak)``), without this extra
    per-neuron state. Uses the fast-sigmoid surrogate gradient by default, matching the
    rest of the package.

    Call convention mirrors ``snn.Leaky`` but carries the extra adaptation state:
    ``spk, mem, b = neuron(input, mem, b)``; ``mem, b = neuron.init_alif(shape)``.
    """

    def __init__(self, beta, tau_a, dt=1.0, gamma=1.6, threshold=1.0, spike_grad=None):
        super().__init__()
        import math
        self.beta = float(beta)
        self.rho = math.exp(-float(dt) / float(tau_a))
        self.tau_a = float(tau_a)
        self.dt = float(dt)
        self.gamma = float(gamma)
        self.threshold = float(threshold)
        if spike_grad is None:
            from snntorch import surrogate
            spike_grad = surrogate.fast_sigmoid(slope=25)
        self.spike_grad = spike_grad

    @staticmethod
    def init_alif(shape):
        """Zero membrane + zero adaptation state of the given ``shape``."""
        return torch.zeros(*shape), torch.zeros(*shape)

    def forward(self, input_, mem, b):
        mem = self.beta * mem + input_                   # leak + integrate this step
        thr = self.threshold + self.gamma * b            # adaptive threshold (prev b)
        spk = self.spike_grad(mem - thr)                 # surrogate-gradient spike
        mem = mem - spk * thr                            # subtract-reset
        b = self.rho * b + spk                           # slow adaptation update
        return spk, mem, b


def _build_homeostasis(memristive_config):
    """Construct the homeostasis module selected by config.

    Returns ``None`` when homeostasis is disabled. By default the
    firing-rate-dependent :class:`HomeostaticRegulariser` (manuscript Eq. homeo)
    is used; setting ``homeostasis_mode='legacy'`` selects the binary
    :class:`HomeostasisDropout` retained for backward compatibility.
    """
    if not memristive_config.get("homeostasis_dropout", False):
        return None
    if memristive_config.get("homeostasis_mode", "regulariser") == "legacy":
        return HomeostasisDropout()
    return HomeostaticRegulariser(
        r_th=memristive_config.get("homeostasis_r_th", 0.8),
        n=memristive_config.get("homeostasis_n", 1.0),
    )


class LayerBuilder:
    """Helper class to build layers based on the configuration."""

    @staticmethod
    def build_fc_layer(in_features, out_features, memristive_config):
        """Builds either a regular Linear layer or MemristiveLinearLayer."""
        if memristive_config.get("ideal", False):
            return nn.Linear(in_features, out_features)
        return MemristiveLinearLayer(in_features, out_features, memristive_config)

    @staticmethod
    def build_conv_layer(
        in_channels, out_channels, kernel_size, stride, padding, memristive_config
    ):
        """Builds either a regular Conv2d layer or MemristiveConv2d."""
        if memristive_config.get("ideal", False):
            return nn.Conv2d(
                in_channels, out_channels, kernel_size, stride=stride, padding=padding
            )
        return MemristiveConv2d(
            in_channels, out_channels, kernel_size, stride, padding, memristive_config
        )


class BaseSNN(nn.Module):
    """Base class to handle the basic operations of an SNN."""

    def __init__(self, beta, spike_grad, num_steps, memristive_config):
        super().__init__()
        self.beta = beta
        self.spike_grad = spike_grad
        self.num_steps = num_steps
        self.memristive_config = memristive_config
        self.homeostasis_threshold = memristive_config.get("homeostasis_threshold", 100)

    def initialize_lif(self, num_layers):
        """Initialize the LIF neurons.

        If the config supplies ``lif_tau_leak`` (a measured device retention) the
        neurons are :class:`DeviceLeaky` -- their membrane decay is the device's own
        trap-discharge relaxation (``beta = exp(-dt/tau_leak)``), the same primitive
        as the RL eligibility trace. Otherwise a plain ``snn.Leaky`` with the given
        ``beta`` is used (unchanged behaviour for image tasks)."""
        tau_leak = self.memristive_config.get("lif_tau_leak")
        if tau_leak is not None:
            dt = self.memristive_config.get("lif_dt", 1.0)
            return [
                DeviceLeaky(tau_leak=tau_leak, dt=dt, spike_grad=self.spike_grad)
                for _ in range(num_layers)
            ]
        return [
            snn.Leaky(beta=self.beta, spike_grad=self.spike_grad)
            for _ in range(num_layers)
        ]


class MSNN(BaseSNN):
    def __init__(
        self, num_inputs, num_hidden, num_outputs, num_steps, beta, memristive_config
    ):
        super().__init__(
            beta,
            spike_grad=None,
            num_steps=num_steps,
            memristive_config=memristive_config,
        )

        # Build layers
        self.fc1 = LayerBuilder.build_fc_layer(
            num_inputs, num_hidden, memristive_config
        )
        self.fc2 = LayerBuilder.build_fc_layer(
            num_hidden, num_outputs, memristive_config
        )

        # Initialize LIF neurons
        self.lif1, self.lif2 = self.initialize_lif(2)

        # Initialize homeostasis module if needed (firing-rate-dependent
        # regulariser by default; legacy binary dropout if configured).
        self.drop_layer = _build_homeostasis(memristive_config)

    def forward(self, x):
        mem1, mem2 = self.lif1.init_leaky(), self.lif2.init_leaky()
        spk1_rec, mem1_rec = [], []
        spk2_rec, mem2_rec = [], []

        # The input is constant current-injection: x does not change across timesteps,
        # and the device-fixed fault mask + frozen PF residual are sampled once and cached
        # (see MemristiveLinearLayer). So fc1(x) is loop-invariant -- computing it ONCE
        # instead of num_steps times is numerically identical and removes the dominant
        # redundant cost (the big input-layer Poole-Frenkel per-synapse tensor, evaluated
        # ~num_steps fewer times). fc2 stays in the loop: it depends on spk1, which varies.
        cur1 = self.fc1(x)

        for _ in range(self.num_steps):
            # First layer (current precomputed above; LIF state still steps per timestep)
            spk1, mem1 = self.lif1(cur1, mem1)

            # Collect spikes
            spk1_rec.append(spk1)
            mem1_rec.append(mem1)

            # Apply drop layer
            if self.drop_layer and len(spk1_rec) >= self.homeostasis_threshold:
                spk1_window = torch.stack(
                    spk1_rec[-self.homeostasis_threshold :], dim=0
                )
                spk1 = self.drop_layer(spk1_window)

            # Second layer
            cur2 = self.fc2(spk1)
            spk2, mem2 = self.lif2(cur2, mem2)

            # Collect spikes
            spk2_rec.append(spk2)
            mem2_rec.append(mem2)

        return torch.stack(spk2_rec, dim=0), torch.stack(mem2_rec, dim=0)


class TemporalMSNN(BaseSNN):
    """Fully-connected spiking network that STREAMS a time-series input across its
    timesteps, so the LIF neurons perform genuine temporal integration of the signal
    as it unfolds -- unlike MSNN, which feeds a static (flattened) input identically
    at every timestep. Built for time-series such as EEG: at network-timestep ``t`` the
    layer sees the ``t``-th sample of the input (the per-channel vector at that time),
    so spike timing and membrane dynamics encode the signal's temporal/spectral
    structure. This is the SNN capability a flattened static encoding discards, and it
    ties to the retention study: the device's slow relaxation integrates the
    streamed signal over time.

    Input ``x`` has shape ``(T, batch, num_inputs)`` where T is the number of input
    time samples; the network runs for ``T`` steps (``num_steps`` is ignored / set to
    T). ``num_inputs`` is the per-timestep feature count (e.g. EEG channels)."""

    def __init__(self, num_inputs, num_hidden, num_outputs, beta, memristive_config):
        super().__init__(
            beta, spike_grad=None, num_steps=0, memristive_config=memristive_config
        )
        self.fc1 = LayerBuilder.build_fc_layer(num_inputs, num_hidden, memristive_config)
        self.fc2 = LayerBuilder.build_fc_layer(num_hidden, num_outputs, memristive_config)
        self.lif1, self.lif2 = self.initialize_lif(2)
        self.drop_layer = _build_homeostasis(memristive_config)

    def forward(self, x):
        # x: (T, batch, num_inputs) -- a streamed time series.
        T = x.shape[0]
        mem1, mem2 = self.lif1.init_leaky(), self.lif2.init_leaky()
        spk1_rec, spk2_rec, mem2_rec = [], [], []

        for t in range(T):
            # fc1 runs on THIS timestep's input slice (not a static vector) -- the LIF
            # state carries history across t, integrating the signal over time.
            cur1 = self.fc1(x[t])
            spk1, mem1 = self.lif1(cur1, mem1)
            spk1_rec.append(spk1)

            if self.drop_layer and len(spk1_rec) >= self.homeostasis_threshold:
                spk1_window = torch.stack(spk1_rec[-self.homeostasis_threshold:], dim=0)
                spk1 = self.drop_layer(spk1_window)

            cur2 = self.fc2(spk1)
            spk2, mem2 = self.lif2(cur2, mem2)
            spk2_rec.append(spk2)
            mem2_rec.append(mem2)

        return torch.stack(spk2_rec, dim=0), torch.stack(mem2_rec, dim=0)


class MCSNN(BaseSNN):
    def __init__(
        self,
        beta,
        spike_grad,
        num_steps,
        batch_size,
        num_kernels,
        num_conv1,
        num_conv2,
        max_pooling,
        num_outputs,
        memristive_config,
        input_shape=(1, 28, 28),  # Default to MNIST input size
        stride=1,
        padding=0,
    ):
        super().__init__(beta, spike_grad, num_steps, memristive_config)

        # Store configurations
        self.batch_size = batch_size
        self.stride = stride
        self.num_conv1 = num_conv1
        self.num_conv2 = num_conv2
        self.max_pooling = max_pooling

        # Build convolutional layers. With ``conv_ideal=True`` the conv layers are
        # standard (ideal) and only the final dense layer carries the Poole-Frenkel
        # device model -- the Joksas et al. methodology (PF on dense layers only),
        # which is fast and avoids the conv PF cost entirely. Homeostasis still acts on
        # the inter-layer spike rates regardless. Default (False) applies PF to the
        # conv layers too (beyond Joksas).
        conv_config = memristive_config
        if memristive_config.get("conv_ideal", False):
            conv_config = {**memristive_config, "ideal": True}
        self.conv1 = LayerBuilder.build_conv_layer(
            input_shape[0], num_conv1, num_kernels, stride, padding, conv_config
        )
        self.conv2 = LayerBuilder.build_conv_layer(
            num_conv1, num_conv2, num_kernels, stride, padding, conv_config
        )

        # Initialize LIF neurons
        self.lif1, self.lif2, self.lif3 = self.initialize_lif(3)

        # Calculate flattened size after convolutions and pooling
        self.num_hidden = self._calculate_hidden_size(
            input_shape, num_kernels, stride, padding, max_pooling
        )

        # Build fully connected layer
        self.fc1 = LayerBuilder.build_fc_layer(
            self.num_hidden, num_outputs, memristive_config
        )

        # Initialize homeostasis modules for each LIF layer before the output
        # (firing-rate-dependent regulariser by default; legacy binary dropout
        # if configured).
        self.homeostasis_dropout1 = _build_homeostasis(memristive_config)
        self.homeostasis_dropout2 = _build_homeostasis(memristive_config)

    def _calculate_hidden_size(
        self, input_shape, kernel_size, stride, padding, max_pooling
    ):
        channels, height, width = input_shape
        # Convolutional layer size calculation (with max pooling) for two convolutional layers
        height = ((height + 2 * padding - kernel_size) // stride + 1) // max_pooling
        width = ((width + 2 * padding - kernel_size) // stride + 1) // max_pooling
        height = ((height + 2 * padding - kernel_size) // stride + 1) // max_pooling
        width = ((width + 2 * padding - kernel_size) // stride + 1) // max_pooling
        return height * width * self.num_conv2

    def forward(self, x):
        mem1, mem2, mem3 = (
            self.lif1.init_leaky(),
            self.lif2.init_leaky(),
            self.lif3.init_leaky(),
        )
        spk1_rec, mem1_rec = [], []
        spk2_rec, mem2_rec = [], []
        spk3_rec, mem3_rec = [], []

        for _ in range(self.num_steps):
            # First convolutional layer + LIF
            cur1 = F.max_pool2d(self.conv1(x), self.max_pooling)
            spk1, mem1 = self.lif1(cur1, mem1)

            # Apply homeostasis dropout after first LIF
            if (
                self.homeostasis_dropout1
                and len(spk1_rec) >= self.homeostasis_threshold
            ):
                spk1_window = torch.stack(
                    spk1_rec[-self.homeostasis_threshold :], dim=0
                )
                spk1 = self.homeostasis_dropout1(spk1_window)

            spk1_rec.append(spk1)
            mem1_rec.append(mem1)

            # Second convolutional layer + LIF
            cur2 = F.max_pool2d(self.conv2(spk1), self.max_pooling)
            spk2, mem2 = self.lif2(cur2, mem2)

            # Apply homeostasis dropout after second LIF
            if (
                self.homeostasis_dropout2
                and len(spk2_rec) >= self.homeostasis_threshold
            ):
                spk2_window = torch.stack(
                    spk2_rec[-self.homeostasis_threshold :], dim=0
                )
                spk2 = self.homeostasis_dropout2(spk2_window)

            spk2_rec.append(spk2)
            mem2_rec.append(mem2)

            # Fully connected layer + LIF. Use the runtime batch dim (not the
            # hardcoded self.batch_size, which breaks on a final partial batch) and
            # reshape rather than view (the PF conv output may be non-contiguous).
            cur3 = self.fc1(spk2.reshape(spk2.shape[0], -1))
            spk3, mem3 = self.lif3(cur3, mem3)
            spk3_rec.append(spk3)
            mem3_rec.append(mem3)

        return torch.stack(spk3_rec), torch.stack(mem3_rec), torch.stack(spk2_rec)


class TemporalMCSNN(MCSNN):
    """Spatiotemporal spiking net: convolution PER FRAME + temporal integration ACROSS
    frames. Combines the two device-grounded pieces -- MemristiveConv2d (spatial features
    per frame, with the binary-spike fast path) and a DeviceLeaky membrane whose leak is
    the measured tau_leak (temporal memory across frames). Where MCSNN feeds a single
    static frame identically at every timestep, this streams a clip: at frame t the conv
    stack runs on x[t], and the LIF membranes carry state across t, so the network
    integrates motion/temporal structure. This is the architecture for spatiotemporal
    data (moving digits, event-camera streams) where BOTH spatial conv and temporal
    memory are needed. Identical construction to MCSNN (same conv dims, hidden-size calc,
    device grounding); only the forward streams over the clip.

    Input ``x`` has shape ``(T, batch, C, H, W)`` -- a clip of T frames. ``num_steps`` is
    ignored; the network runs for T frames."""

    def forward(self, x):
        # x: (T, batch, C, H, W)
        T = x.shape[0]
        mem1, mem2, mem3 = (
            self.lif1.init_leaky(),
            self.lif2.init_leaky(),
            self.lif3.init_leaky(),
        )
        spk1_rec, spk2_rec, spk3_rec, mem3_rec = [], [], [], []

        for t in range(T):
            # conv stack runs on THIS frame; LIF state carries across frames.
            cur1 = F.max_pool2d(self.conv1(x[t]), self.max_pooling)
            spk1, mem1 = self.lif1(cur1, mem1)
            if (self.homeostasis_dropout1
                    and len(spk1_rec) >= self.homeostasis_threshold):
                spk1 = self.homeostasis_dropout1(
                    torch.stack(spk1_rec[-self.homeostasis_threshold:], dim=0))
            spk1_rec.append(spk1)

            cur2 = F.max_pool2d(self.conv2(spk1), self.max_pooling)
            spk2, mem2 = self.lif2(cur2, mem2)
            if (self.homeostasis_dropout2
                    and len(spk2_rec) >= self.homeostasis_threshold):
                spk2 = self.homeostasis_dropout2(
                    torch.stack(spk2_rec[-self.homeostasis_threshold:], dim=0))
            spk2_rec.append(spk2)

            cur3 = self.fc1(spk2.reshape(spk2.shape[0], -1))
            spk3, mem3 = self.lif3(cur3, mem3)
            spk3_rec.append(spk3)
            mem3_rec.append(mem3)

        return torch.stack(spk3_rec), torch.stack(mem3_rec), torch.stack(spk2_rec)
\end{lstlisting}

\begin{lstlisting}[caption={mnn\_torch/data.py},label={lst:mnn-src-data}]
"""Dataset loading and preparation helpers for the mnn-torch experiments.

Two dataset paths feed the memristive spiking-network studies:

* **MNIST** -- the fault-recovery / homeostasis sweeps.
  :func:`ensure_mnist` downloads it once (into :func:`mnn_torch.paths.data_dir`
  by default) so that pooled worker processes can open it with ``download=False``;
  :func:`mnist_loaders` builds seeded train/test ``DataLoader`` pairs, optionally
  subsetting for a fast smoke run.

* **BCI Competition IV-2a** (BNCI2014_001, 4-class motor imagery) -- the
  biosignal-classification study. :func:`prep_bci2a` downloads and
  epochs it via MOABB/MNE and writes a shippable ``.npz`` so the heavy EEG
  dependency tree never has to enter the training container.

``torch``/``torchvision`` and ``moabb``/``mne`` are imported *lazily* inside the
functions so that importing :mod:`mnn_torch.data` (and hence :mod:`mnn_torch`)
never drags in those heavy trees; only the caller who actually needs a loader
pays for the import.
"""
from __future__ import annotations

import os

import numpy as np

from . import paths

__all__ = [
    "mnist_loaders",
    "ensure_mnist",
    "prep_bci2a",
    # BCI IV-2a preprocessing constants (module-level so callers/tests can read them)
    "FMIN",
    "FMAX",
    "TMIN",
    "TMAX",
    "N_EEG",
    "DECIMATE",
]


# --------------------------------------------------------------------------
# MNIST (fault-recovery / homeostasis sweeps)
# --------------------------------------------------------------------------
def ensure_mnist(data_dir=None):
    """Download MNIST once (train + test) so workers can use ``download=False``.

    ``data_dir`` defaults to :func:`mnn_torch.paths.data_dir`. Returns the
    resolved data directory as a string so callers can thread it to workers.
    """
    import torchvision
    import torchvision.transforms as T

    data_dir = str(paths.data_dir()) if data_dir is None else str(data_dir)
    tfm = T.Compose([T.ToTensor()])
    torchvision.datasets.MNIST(data_dir, train=True, download=True, transform=tfm)
    torchvision.datasets.MNIST(data_dir, train=False, download=True, transform=tfm)
    return data_dir


def mnist_loaders(data_dir=None, seed=0, batch_size=64, train_subset=None,
                  test_subset=None):
    """Build seeded MNIST train/test ``DataLoader`` pairs.

    Lifted from ``gate_homeostasis_recovery._data_loaders``: the train loader is
    shuffled under a per-run ``torch.Generator`` seed and drops the last partial
    batch; the test loader is deterministic. ``train_subset`` / ``test_subset``
    (if given) take the leading ``n`` examples for a fast smoke run. Assumes
    MNIST is already present (see :func:`ensure_mnist`); opens with
    ``download=False``.
    """
    import torch
    import torchvision
    import torchvision.transforms as T
    from torch.utils.data import DataLoader, Subset

    data_dir = str(paths.data_dir()) if data_dir is None else str(data_dir)
    tfm = T.Compose([T.ToTensor(), T.Normalize((0,), (1,))])
    train = torchvision.datasets.MNIST(data_dir, train=True, download=False, transform=tfm)
    test = torchvision.datasets.MNIST(data_dir, train=False, download=False, transform=tfm)
    if train_subset:
        train = Subset(train, range(train_subset))
    if test_subset:
        test = Subset(test, range(test_subset))
    g = torch.Generator().manual_seed(seed)
    train_loader = DataLoader(train, batch_size=batch_size, shuffle=True,
                              drop_last=True, generator=g)
    test_loader = DataLoader(test, batch_size=batch_size, shuffle=False,
                             drop_last=True)
    return train_loader, test_loader


# --------------------------------------------------------------------------
# BCI Competition IV-2a (biosignal classification)
# --------------------------------------------------------------------------
# Standard 4-class motor-imagery preprocessing (see prep_bci2a docstring).
FMIN, FMAX = 8.0, 30.0          # mu + beta motor-imagery band
TMIN, TMAX = 0.5, 2.5           # seconds post-cue (avoid the cue transient)
N_EEG = 22                      # drop the 3 EOG channels
DECIMATE = 4                    # 250 Hz -> 62.5 Hz: the band is <=30 Hz so this is
                                # safely above Nyquist, and cuts the flatten dim from
                                # 22x501=11022 to 22x125=2750 (MNIST-scale, no bottleneck).


def _load_bci2a_epochs():
    """Download (cached by MOABB) and epoch BCI IV-2a. Returns X[n,ch,t], y[n] int 0-3,
    subj[n] int, and the class-name list. MOABB/MNE imported lazily."""
    from moabb.datasets import BNCI2014_001
    from moabb.paradigms import MotorImagery

    # 4-class MI, our band-pass and window; MOABB returns z-scoring-free volts.
    paradigm = MotorImagery(n_classes=4, fmin=FMIN, fmax=FMAX, tmin=TMIN, tmax=TMAX,
                            channels=None, resample=None)
    ds = BNCI2014_001()
    X, labels, meta = paradigm.get_data(dataset=ds)        # X: (n, n_chan, n_time)
    X = X[:, :N_EEG, :].astype(np.float32)                 # EEG channels only
    if DECIMATE > 1:                                       # anti-alias-safe: band <= 30 Hz
        X = X[:, :, ::DECIMATE]                            # 501 -> ~125 timepoints

    classes = sorted(np.unique(labels).tolist())           # e.g. ['feet','left_hand',...]
    cls_to_int = {c: i for i, c in enumerate(classes)}
    y = np.array([cls_to_int[l] for l in labels], dtype=np.int64)
    subj = meta["subject"].to_numpy().astype(np.int64)

    # Per-trial, per-channel z-score (the SNN front-end expects standardised input, and
    # this matches the biosignal.py convention in the sibling repo).
    mu = X.mean(axis=2, keepdims=True)
    sd = X.std(axis=2, keepdims=True) + 1e-7
    X = (X - mu) / sd
    return X, y, subj, classes


def _csp_lda_sanity(X, y, subj, seed=0):
    """Within-subject CSP+LDA 5-fold accuracy, averaged over subjects -- the clean-ceiling
    gate. Uses MNE's CSP on the covariance of the band-passed epochs."""
    from mne.decoding import CSP
    from sklearn.discriminant_analysis import LinearDiscriminantAnalysis
    from sklearn.pipeline import Pipeline
    from sklearn.model_selection import StratifiedKFold, cross_val_score

    accs = []
    for s in np.unique(subj):
        m = subj == s
        Xs, ys = X[m], y[m]
        clf = Pipeline([("csp", CSP(n_components=8, reg="ledoit_wolf", log=True)),
                        ("lda", LinearDiscriminantAnalysis())])
        skf = StratifiedKFold(5, shuffle=True, random_state=seed)
        a = cross_val_score(clf, Xs, ys, cv=skf, scoring="accuracy").mean()
        accs.append(a)
        print(f"  subject {s:2d}: CSP+LDA 4-class acc = {a:.3f}  (n={m.sum()})")
    return float(np.mean(accs)), float(np.std(accs))


def prep_bci2a(out=None, skip_sanity=False):
    """Pre-cache the BCI Competition IV-2a (BNCI2014_001) 4-class motor-imagery EEG
    as a shippable ``.npz``.

    Downloads + epochs the dataset via MOABB, optionally runs the within-subject
    CSP+LDA clean-ceiling sanity decode, and writes ``X/y/subj/classes`` plus the
    band-pass/window metadata to ``out`` (default:
    ``paths.data_dir()/"bci2a_epochs.npz"``). Returns the absolute output path.

    MOABB/MNE are heavy and would re-download on every batch run, so the download +
    epoching happens ONCE here (in an environment that already has moabb/mne) and
    only the dense array is shipped -- exactly as the SiOx ``.mat`` is shipped.
    """
    out = paths.data_dir() / "bci2a_epochs.npz" if out is None else out
    out = os.path.abspath(str(out))

    print("[prep] downloading + epoching BCI IV-2a (BNCI2014_001) via MOABB ...")
    X, y, subj, classes = _load_bci2a_epochs()
    print(f"[prep] epochs: X={X.shape} (trials, ch, time)  classes={classes}  "
          f"subjects={np.unique(subj).tolist()}")
    print(f"[prep] class balance: {np.bincount(y).tolist()}")

    mean_acc = std_acc = float("nan")
    if not skip_sanity:
        print("[prep] CSP+LDA within-subject sanity decode (clean-ceiling gate) ...")
        mean_acc, std_acc = _csp_lda_sanity(X, y, subj)
        print(f"[prep] CSP+LDA mean 4-class accuracy = {mean_acc:.3f} +/- {std_acc:.3f} "
              f"(chance 0.25; published CSP+LDA ~0.67)")
        if mean_acc < 0.40:
            print("[prep] WARNING: clean ceiling is low -- check band-pass/window before "
                  "spending SNN effort.")

    os.makedirs(os.path.dirname(out), exist_ok=True)
    np.savez_compressed(out, X=X, y=y, subj=subj, classes=np.array(classes),
                        fmin=FMIN, fmax=FMAX, tmin=TMIN, tmax=TMAX,
                        csp_lda_acc=mean_acc, csp_lda_std=std_acc,
                        sfreq=250.0, n_eeg=N_EEG)
    print(f"[prep] wrote {out}  ({os.path.getsize(out)/1e6:.1f} MB)")
    return out


def main(argv=None):
    """CLI wrapper for :func:`prep_bci2a` (mirrors the old prep_bci2a.py entry).

    ``python -m mnn_torch.data [--out PATH] [--skip-sanity]``
    """
    import argparse

    ap = argparse.ArgumentParser(description="Pre-cache BCI IV-2a epochs as a .npz")
    ap.add_argument("--out", default=None,
                    help="output .npz path (default: paths.data_dir()/bci2a_epochs.npz)")
    ap.add_argument("--skip-sanity", action="store_true")
    a = ap.parse_args(argv)
    prep_bci2a(out=a.out, skip_sanity=a.skip_sanity)


if __name__ == "__main__":
    main()
\end{lstlisting}

\begin{lstlisting}[caption={mnn\_torch/training.py},label={lst:mnn-src-training}]
"""Training, config, and fault-recovery sweep logic for the memristive SNN studies.

This module owns the *compute* for the publication-scale homeostatic
fault-recovery study and its two diagnostic slices, woven beside
:mod:`mnn_torch.models` / :mod:`mnn_torch.effects` (which own the device
primitives). Nothing here plots inside a core and nothing writes files inside a
core: the reusable pieces (:func:`make_config`, :func:`train_snn`,
:func:`evaluate_snn`, :func:`run_condition`, :func:`gate_homeostasis_sweep`,
:func:`dose_response`, :func:`isolate_collapse`, :func:`precompute_device_params`)
RETURN plain data. Result I/O and the heavy full-scale drivers live in
:func:`main` (argparse ``--gate/--dose/--isolate`` x ``--quick/--full``) and are
routed through :mod:`mnn_torch.paths`, so notebooks and ``python -m
mnn_torch.training`` agree on where grids land regardless of CWD.

``torch``/``torchvision`` are imported *lazily* inside the functions so that
importing :mod:`mnn_torch.training` never drags in the heavy tree.

Science (unchanged from the experiments/ scripts it supersedes)
---------------------------------------------------------------
A fully connected memristive spiking net (MSNN, 784->100->10) is trained on
MNIST with the measured SiOx Poole-Frenkel device prior (Eq. forward-g). We
sweep the device-failure probability ``p`` and the stuck-fault polarity:

  * stuck-on  : saturated/shorted synapses drive their postsynaptic line
                HYPERACTIVE -> exactly what the firing-rate homeostatic
                regulariser (Eq. homeo) suppresses -> recovery SHOULD be large.
  * stuck-off : dead/open synapses go silent -> the homeostat has nothing to act
                on -> the mechanism-NEGATIVE control (little/no recovery).

For every (polarity, p) cell we train two arms -- homeostasis on vs off -- and
report the recovery gap (homeo - no_homeo). ``p = 0`` (faults disabled) is the
clean ceiling and its "recovery" is the generic-regularisation baseline the
fault-driven recovery must rise above. The mechanism test contrasts peak
stuck-on recovery against stuck-off, with bootstrap CIs across seeds.

The production entry (spawn Pool, dataset matrix, MCSNN path) folds in behind
``--full``; the default ``--quick`` runs a tiny in-kernel-safe sweep.
"""
from __future__ import annotations

import os
import subprocess
import warnings
from concurrent.futures import ProcessPoolExecutor, as_completed

import numpy as np

from . import paths

warnings.filterwarnings("ignore")

__all__ = [
    "make_config",
    "_make_config",
    "evaluate_snn",
    "train_snn",
    "run_condition",
    "precompute_device_params",
    "gate_homeostasis_sweep",
    "dose_response",
    "isolate_collapse",
    "main",
]


# ==========================================================================
# Config builders
#
# The three source scripts (gate_homeostasis_recovery, dose_response,
# isolate_collapse) each had a near-identical memristive-config builder. They
# are unified here into one ``make_config`` that PRESERVES every key and value
# from all three, with the union of their optional knobs exposed as keyword
# arguments defaulted to each script's original value:
#
#   * ``prob`` -> disturbance_probability (all three).
#   * ``homeo`` -> homeostasis_dropout (all three).
#   * ``nonlinear`` -> PF (True) vs Ohmic (False) forward map. gate/dose used
#     True; isolate toggled it. Default True.
#   * ``disturb_mode`` -> "device_fixed" (gate default), "device" (dose/isolate),
#     or "fixed" (isolate). Default "device_fixed".
#   * ``stuck_polarity`` -> "on"/"off"/"split". gate swept it; dose/isolate never
#     set it (so it stayed the layer default). Default None => key omitted, which
#     reproduces dose/isolate exactly.
#   * ``disturb_conductance``: gate used ``prob > 0`` (p=0 truly disables the
#     path); dose used the explicit ``faulty`` flag; isolate hardwired True.
#     Default None => derive as ``prob > 0`` (the gate rule); pass an explicit
#     bool to reproduce dose (faulty) / isolate (True).
#   * ``homeostasis_r_th`` / ``homeostasis_n`` / ``homeostasis_threshold``:
#     identical across all three (0.8 / 1.0 / 5); kept as defaults.
#   * ``pf_frozen_variability``: gate set it True (fixed-device reading); dose/
#     isolate omitted it (layer default). Default None => omit unless requested.
# ==========================================================================
def make_config(dev_params, *, prob, homeo, nonlinear=True,
                disturb_mode="device_fixed", stuck_polarity=None,
                disturb_conductance=None, pf_frozen_variability=None,
                homeostasis_r_th=0.8, homeostasis_n=1.0, homeostasis_threshold=5):
    """Build a memristive config for a single training run.

    ``dev_params`` is the 8-tuple from :func:`precompute_device_params`
    ``(G_off, G_on, R, c, d_epsilon, pf_slopes, pf_intercepts, pf_covariance)``.

    See the module comment for how every key from the three original builders is
    preserved. When ``disturb_conductance is None`` the disturbance path is
    enabled iff ``prob > 0`` (the gate rule: ``prob == 0`` is a true clean
    ceiling, not the fault sampler run at probability 0).
    """
    G_off, G_on, R, c, d_eps, slopes, intercepts, cov = dev_params
    if disturb_conductance is None:
        disturb_conductance = prob > 0.0
    cfg = {
        "ideal": False,
        "k_V": 0.5,
        "G_off": float(G_off),
        "G_on": float(G_on),
        "R": R,
        "c": c,
        "d_epsilon": d_eps,
        # Poole-Frenkel nonlinear forward map (Eq. forward-g).
        "nonlinear": nonlinear,
        "pf_slopes": slopes,
        "pf_intercepts": intercepts,
        "pf_covariance": cov,
        "pf_stochastic": True,
        # faults (swept).
        "disturb_conductance": disturb_conductance,
        "disturb_mode": disturb_mode,
        "disturbance_probability": float(prob),
        "disturbance_sigma_g": 0.5,
        # homeostasis (sigmoidal firing-rate regulariser, Eq. homeo).
        "homeostasis_dropout": homeo,
        "homeostasis_mode": "regulariser",
        "homeostasis_r_th": homeostasis_r_th,
        "homeostasis_n": homeostasis_n,
        "homeostasis_threshold": homeostasis_threshold,
    }
    # Optional keys only added when set, so dose/isolate configs are reproduced
    # byte-for-byte (they never emitted these keys).
    if stuck_polarity is not None:
        cfg["disturbance_stuck_polarity"] = stuck_polarity
    if pf_frozen_variability is not None:
        cfg["pf_frozen_variability"] = pf_frozen_variability
    return cfg


def _make_config(dev_params, *, prob, homeo, disturb_mode="device_fixed",
                 stuck_polarity="split"):
    """Gate-sweep config builder (compat alias, exact keys of the original
    ``gate_homeostasis_recovery._make_config``).

    Sets ``pf_frozen_variability=True`` (the fixed-device reading, frozen PF
    residual) and threads ``stuck_polarity`` through -- i.e. the gate arm.
    """
    return make_config(dev_params, prob=prob, homeo=homeo, nonlinear=True,
                       disturb_mode=disturb_mode, stuck_polarity=stuck_polarity,
                       disturb_conductance=(prob > 0.0),
                       pf_frozen_variability=True)


# ==========================================================================
# Evaluate / train (shared by every arm)
# ==========================================================================
def evaluate_snn(net, loader, device):
    """Rate-coded top-1 accuracy (%) of ``net`` on ``loader``.

    Prediction = argmax of the summed output spike train (rate code), matching
    the training loss. Flattens each input to ``(B, features)`` for the MSNN path.
    """
    import torch

    net.eval()
    correct = total = 0
    with torch.no_grad():
        for x, y in loader:
            x = x.view(x.shape[0], -1).to(device)
            y = y.to(device)
            spk, _ = net(x)
            pred = spk.sum(dim=0).argmax(dim=1)
            correct += (pred == y).sum().item()
            total += y.shape[0]
    return 100.0 * correct / total


def train_snn(cfg, loaders, device, *, seed, epochs, num_steps, lr):
    """Train an MSNN(784->100->10) under ``cfg`` and return ``(final_acc, hist)``.

    ``hist`` is the per-epoch test accuracy list (as dose/isolate returned).
    Rate-coded cross-entropy on the summed output spikes; Adam optimiser.
    """
    import torch
    import torch.nn as nn

    from .models import MSNN

    torch.manual_seed(seed)
    np.random.seed(seed)
    train_loader, test_loader = loaders
    net = MSNN(28 * 28, 100, 10, num_steps=num_steps, beta=0.95,
               memristive_config=cfg).to(device)
    opt = torch.optim.Adam(net.parameters(), lr=lr)
    loss_fn = nn.CrossEntropyLoss()
    hist = []
    for _ in range(epochs):
        net.train()
        for x, y in train_loader:
            x = x.view(x.shape[0], -1).to(device)
            y = y.to(device)
            spk, _ = net(x)
            loss = loss_fn(spk.sum(dim=0), y)  # rate-coded CE on summed spikes
            opt.zero_grad()
            loss.backward()
            opt.step()
        hist.append(evaluate_snn(net, test_loader, device))
    return hist[-1], hist


# ==========================================================================
# Gate-sweep worker (picklable; imports torch lazily)
# ==========================================================================
def run_condition(job):
    """Train + evaluate one gate-sweep cell. Returns the ``job`` dict + ``acc``.

    Picklable entry point for the ProcessPool. ``job`` carries every scalar plus
    the precomputed device-parameter arrays, so no global state crosses the
    process boundary. Lifted verbatim from
    ``gate_homeostasis_recovery.run_condition`` (routed through this module's
    :func:`_make_config`, :func:`mnist_loaders`, :func:`evaluate_snn`).
    """
    import torch
    import torch.nn as nn

    from .data import mnist_loaders
    from .models import MSNN

    if job["threads"]:
        torch.set_num_threads(job["threads"])
    device = torch.device(job["device"])

    torch.manual_seed(job["seed"])
    np.random.seed(job["seed"])

    cfg = _make_config(job["dev_params"], prob=job["prob"], homeo=job["homeo"],
                       disturb_mode=job["disturb_mode"],
                       stuck_polarity=job["stuck_polarity"])
    train_loader, test_loader = mnist_loaders(
        job["data"], seed=job["seed"], batch_size=job["batch_size"],
        train_subset=job["train_subset"], test_subset=job["test_subset"])
    net = MSNN(28 * 28, 100, 10, num_steps=job["num_steps"], beta=0.95,
               memristive_config=cfg).to(device)
    opt = torch.optim.Adam(net.parameters(), lr=job["lr"])
    loss_fn = nn.CrossEntropyLoss()
    for _ in range(job["epochs"]):
        net.train()
        for x, y in train_loader:
            x = x.view(x.shape[0], -1).to(device)
            y = y.to(device)
            spk, _ = net(x)
            loss = loss_fn(spk.sum(dim=0), y)  # rate-coded CE on summed spikes
            opt.zero_grad()
            loss.backward()
            opt.step()
    job = dict(job)
    job["acc"] = evaluate_snn(net, test_loader, device)
    return job


# ==========================================================================
# Device-parameter precompute (fit the SiOx prior ONCE)
# ==========================================================================
def precompute_device_params(data_dir=None):
    """Fit the SiOx device parameters ONCE; return the 8-tuple broadcast to workers.

    Reads the bundled multistate ``.mat`` via :func:`mnn_torch.paths.device_data_mat`
    (``data_dir`` overrides the directory), fits ``G_off/G_on`` and the
    Poole-Frenkel regression, and returns
    ``(G_off, G_on, R, c, d_epsilon, pf_slopes, pf_intercepts, pf_covariance)``.
    """
    from .devices import load_SiOx_multistate
    from .effects import (
        compute_PooleFrenkel_parameters,
        compute_PooleFrenkel_regression_parameters,
    )

    if data_dir is None:
        mat = str(paths.device_data_mat())
    else:
        mat = os.path.join(str(data_dir), "SiO_x-multistate-data.mat")
    exp = load_SiOx_multistate(mat)
    G_off, G_on, R, c, d_eps = compute_PooleFrenkel_parameters(exp)
    slopes, intercepts, cov = compute_PooleFrenkel_regression_parameters(R, c, d_eps)
    return (G_off, G_on, R, c, d_eps, slopes, intercepts, cov)


# ==========================================================================
# Small helpers (git provenance, device pick, bootstrap CI)
# ==========================================================================
def git_commit():
    try:
        return subprocess.check_output(
            ["git", "rev-parse", "--short", "HEAD"],
            cwd=str(paths.data_dir().parent),
            stderr=subprocess.DEVNULL,
        ).decode().strip()
    except Exception:
        return "unknown"


def pick_device(requested):
    import torch

    if requested != "auto":
        return requested
    if torch.backends.mps.is_available():
        return "mps"
    if torch.cuda.is_available():
        return "cuda"
    return "cpu"


def _boot_ci(vals, n_boot=10000, seed=0):
    """95% bootstrap CI of the mean (percentile method). Lifted from
    entry_gate_sweep._boot_ci."""
    vals = np.asarray(vals, dtype=float)
    if len(vals) < 2:
        v = float(vals.mean()) if len(vals) else float("nan")
        return v, v
    rng = np.random.default_rng(seed)
    boot = [vals[rng.integers(0, len(vals), len(vals))].mean() for _ in range(n_boot)]
    return float(np.percentile(boot, 2.5)), float(np.percentile(boot, 97.5))


# ==========================================================================
# Gate-sweep CORE (serial, returns plain data -- for the in-kernel quick run)
# ==========================================================================
def gate_homeostasis_sweep(*, dev_params=None, data_dir=None, seeds=(0,),
                           probs=(0.0, 0.1, 0.2, 0.3, 0.4, 0.5),
                           polarities=("on", "off"), disturb_mode="device_fixed",
                           epochs=4, num_steps=8, train_subset=6000,
                           test_subset=2000, batch_size=64, lr=5e-4,
                           device="cpu", threads=0, workers=1, verbose=True):
    """Run the homeostatic fault-recovery sweep and RETURN the result grid dict.

    Serial by default (``workers=1``) so it is safe to call in a notebook kernel;
    ``workers>1`` fans the independent (seed x polarity x prob x arm) cells out
    over a :class:`ProcessPoolExecutor` of :func:`run_condition` workers.

    Returns a dict with the raw accuracy grids ``acc[f"{pol}|{arm}"] -> (n_seeds,
    n_probs)``, the per-cell recovery + bootstrap CIs, the fault-peak per
    polarity, and (when both ``on`` and ``off`` are swept) the mechanism-test
    contrast. No file I/O, no plotting -- :func:`main` handles persistence.
    """
    seeds = list(seeds)
    probs = list(probs)
    polarities = list(polarities)
    if dev_params is None:
        dev_params = precompute_device_params(data_dir)
    data = str(paths.data_dir()) if data_dir is None else str(data_dir)

    # ---- build the independent job list ----
    jobs = []
    for si, seed in enumerate(seeds):
        for pol in polarities:
            for pi, prob in enumerate(probs):
                for homeo in (False, True):
                    jobs.append(dict(
                        si=si, seed=seed, pi=pi, prob=prob, polarity=pol,
                        stuck_polarity=pol, homeo=homeo,
                        disturb_mode=disturb_mode,
                        epochs=epochs, num_steps=num_steps,
                        train_subset=train_subset, test_subset=test_subset,
                        batch_size=batch_size, lr=lr, data=data,
                        device=device, threads=threads,
                        dev_params=dev_params,
                    ))
    if verbose:
        print(f"{len(jobs)} independent conditions "
              f"({len(seeds)} seeds x {len(polarities)} polarities x "
              f"{len(probs)} probs x 2 arms)")

    # results[(polarity, 'homeo'|'no')] -> (n_seeds, n_probs)
    acc = {}
    for pol in polarities:
        acc[(pol, "homeo")] = np.zeros((len(seeds), len(probs)))
        acc[(pol, "no")] = np.zeros((len(seeds), len(probs)))

    def store(done):
        arm = "homeo" if done["homeo"] else "no"
        acc[(done["polarity"], arm)][done["si"], done["pi"]] = done["acc"]
        if verbose:
            print(f"  done: seed{done['seed']} {done['polarity']:5s} "
                  f"p={done['prob']:.2f} {'homeo ' if done['homeo'] else 'nohomeo'} "
                  f"-> {done['acc']:.2f}%")

    if workers > 1:
        with ProcessPoolExecutor(max_workers=workers) as ex:
            futs = [ex.submit(run_condition, j) for j in jobs]
            for fut in as_completed(futs):
                store(fut.result())
    else:
        for j in jobs:
            store(run_condition(j))

    # ---- summarise (recovery, per-cell CIs, fault-peak, mechanism test) ----
    probs_arr = np.asarray(probs, dtype=float)
    fault_j = [j for j, p in enumerate(probs) if p > 0.0]

    summary = {}
    for pol in polarities:
        nh, h = acc[(pol, "no")], acc[(pol, "homeo")]
        rec = h - nh
        cell = []
        for j, p in enumerate(probs):
            lo, hi = _boot_ci(rec[:, j], seed=j)
            cell.append(dict(p=p, no_homeo=float(nh[:, j].mean()),
                             homeo=float(h[:, j].mean()),
                             recovery=float(rec[:, j].mean()), ci_lo=lo, ci_hi=hi))
        peak = None
        if fault_j:
            jpk = max(fault_j, key=lambda j: rec[:, j].mean())
            plo, phi = _boot_ci(rec[:, jpk], seed=jpk)
            peak = dict(p=probs[jpk], recovery=float(rec[:, jpk].mean()),
                        ci_lo=plo, ci_hi=phi)
        base = float((h - nh)[:, 0].mean()) if probs_arr[0] == 0.0 else float("nan")
        summary[pol] = dict(cells=cell, peak=peak, ideal_effect=base)

    mechanism = None
    if "on" in polarities and "off" in polarities and fault_j:
        rec_on = acc[("on", "homeo")] - acc[("on", "no")]
        rec_off = acc[("off", "homeo")] - acc[("off", "no")]
        j_on = max(fault_j, key=lambda j: rec_on[:, j].mean())
        diff = rec_on[:, j_on] - rec_off[:, j_on]   # paired by seed
        dlo, dhi = _boot_ci(diff, seed=99)
        if dlo > 0:
            verdict = "MECHANISM-SPECIFIC (on > off, CI excludes 0)"
        elif dhi < 0:
            verdict = "REVERSED (off > on, CI excludes 0) -- unexpected"
        else:
            verdict = "NOT separable (CI crosses 0; need more seeds)"
        mechanism = dict(p=probs[j_on], diff=float(diff.mean()),
                         ci_lo=dlo, ci_hi=dhi, verdict=verdict)

    return {
        "probs": probs, "polarities": polarities, "seeds": len(seeds),
        "seed_values": seeds, "epochs": epochs, "num_steps": num_steps,
        "disturb_mode": disturb_mode, "train_subset": train_subset,
        "test_subset": test_subset, "dataset": "mnist",
        "acc": {f"{pol}|{arm}": acc[(pol, arm)]
                for pol in polarities for arm in ("no", "homeo")},
        "summary": summary, "mechanism_test": mechanism,
    }


def _print_gate_summary(grid):
    """Human-readable dump of a :func:`gate_homeostasis_sweep` grid."""
    probs = grid["probs"]
    print(f"\n================ SWEEP SUMMARY (mean over {grid['seeds']} seed[s], "
          "95% bootstrap CI on recovery) ================")
    acc = grid["acc"]
    for pol in grid["polarities"]:
        nh = acc[f"{pol}|no"]
        h = acc[f"{pol}|homeo"]
        rec = (h - nh).mean(0)
        s = grid["summary"][pol]
        print(f"\n  --- stuck-{pol} ---")
        print(f"  {'p':>5}  {'no-homeo':>9}  {'homeo':>9}  {'recovery':>9}  {'95% CI':>16}")
        for cell in s["cells"]:
            print(f"  {cell['p']:5.2f}  {cell['no_homeo']:8.2f}%  {cell['homeo']:8.2f}%  "
                  f"{cell['recovery']:+8.2f}   [{cell['ci_lo']:+.2f},{cell['ci_hi']:+.2f}]")
        pk = s["peak"]
        if pk is not None:
            print(f"  ideal effect (p=0) = {s['ideal_effect']:+.2f} pp ; "
                  f"peak fault recovery = {pk['recovery']:+.2f} pp "
                  f"[{pk['ci_lo']:+.2f},{pk['ci_hi']:+.2f}] at p={pk['p']:.2f}")
    mt = grid["mechanism_test"]
    if mt is not None:
        print(f"\n  MECHANISM TEST @ p={mt['p']:.2f}: recovery(on)-recovery(off) = "
              f"{mt['diff']:+.2f} pp [{mt['ci_lo']:+.2f},{mt['ci_hi']:+.2f}] -> {mt['verdict']}")


# ==========================================================================
# Dose-response CORE (single fault rate; returns plain data)
# ==========================================================================
def dose_response(rate, *, dev_params=None, data_dir=None, seed=0, epochs=5,
                  num_steps=5, train_subset=4000, test_subset=1500,
                  batch_size=64, lr=5e-4, device=None, verbose=True):
    """Dose-response slice at ONE device-failure ``rate``; RETURN the result dict.

    ``rate == 0`` produces the fault-rate-independent reference conditions
    (``ideal`` / ``ideal+homeo``); ``rate > 0`` measures ``faulty`` /
    ``faulty+homeo`` and their ``recovery`` gap. Uses ``disturb_mode="device"``
    (the dose-response setting) and the Poole-Frenkel forward map. No file I/O.
    """
    import torch

    if device is None:
        device = torch.device(pick_device("auto"))
    elif isinstance(device, str):
        device = torch.device(device)
    if dev_params is None:
        dev_params = precompute_device_params(data_dir)

    from .data import ensure_mnist, mnist_loaders
    data = ensure_mnist(data_dir)
    loaders = mnist_loaders(data, seed=seed, batch_size=batch_size,
                            train_subset=train_subset, test_subset=test_subset)

    out = {"rate": rate, "device": str(device), "epochs": epochs,
           "num_steps": num_steps, "train_subset": train_subset, "seed": seed}

    if rate == 0.0:
        arms = [("ideal", False), ("ideal+homeo", True)]
        faulty = False
    else:
        arms = [("faulty", False), ("faulty+homeo", True)]
        faulty = True
    for name, homeo in arms:
        cfg = make_config(dev_params, prob=rate, homeo=homeo, nonlinear=True,
                          disturb_mode="device", disturb_conductance=faulty)
        acc, hist = train_snn(cfg, loaders, device, seed=seed, epochs=epochs,
                              num_steps=num_steps, lr=lr)
        out[name] = acc
        out[name + "_hist"] = hist
        if verbose:
            print(f"[rate={rate}] {name:14s} {acc:.2f}%  "
                  f"hist={['%.1f' % a for a in hist]}", flush=True)
    if rate != 0.0:
        out["recovery"] = out["faulty+homeo"] - out["faulty"]
        if verbose:
            print(f"[rate={rate}] recovery = {out['recovery']:+.2f}%", flush=True)
    return out


# ==========================================================================
# Isolate-collapse CORE (one (map, mode) combination; returns plain data)
# ==========================================================================
def isolate_collapse(map, mode, *, dev_params=None, data_dir=None, rate=0.5,
                     seed=0, epochs=5, num_steps=5, train_subset=4000,
                     test_subset=1500, batch_size=64, lr=5e-4, device="cpu",
                     verbose=True):
    """Isolate the faulty-baseline collapse: one ``(map, mode)`` combination.

    ``map`` in {"pf", "ohmic"} toggles the forward map (Poole-Frenkel vs Ohmic
    ``I=V*G``); ``mode`` in {"device", "fixed"} toggles the disturbance sampler.
    Measures ``faulty`` / ``faulty+homeo`` at fixed ``rate`` and their
    ``recovery``. RETURNS the result dict (no file I/O). CPU by default, matching
    the parallel-on-CPU dose sweep.
    """
    import torch

    device = torch.device(device) if isinstance(device, str) else device
    nonlinear = map == "pf"
    if dev_params is None:
        dev_params = precompute_device_params(data_dir)

    from .data import ensure_mnist, mnist_loaders
    data = ensure_mnist(data_dir)
    loaders = mnist_loaders(data, seed=seed, batch_size=batch_size,
                            train_subset=train_subset, test_subset=test_subset)

    out = {"map": map, "mode": mode, "rate": rate, "epochs": epochs,
           "num_steps": num_steps, "train_subset": train_subset, "seed": seed}
    for name, homeo in [("faulty", False), ("faulty+homeo", True)]:
        cfg = make_config(dev_params, prob=rate, homeo=homeo, nonlinear=nonlinear,
                          disturb_mode=mode, disturb_conductance=True)
        acc, hist = train_snn(cfg, loaders, device, seed=seed, epochs=epochs,
                              num_steps=num_steps, lr=lr)
        out[name] = acc
        out[name + "_hist"] = hist
        if verbose:
            print(f"[{map}/{mode}] {name:14s} {acc:.2f}%  "
                  f"hist={['%.1f' % a for a in hist]}", flush=True)
    out["recovery"] = out["faulty+homeo"] - out["faulty"]
    if verbose:
        print(f"[{map}/{mode}] recovery = {out['recovery']:+.2f}%", flush=True)
    return out


# ==========================================================================
# CLI drivers (argparse dispatch; result I/O routed through paths)
# ==========================================================================
def _gate_main(args):
    """Gate-sweep driver: quick in-kernel sweep, or the production spawn-Pool
    dataset-matrix sweep behind ``--full``. Writes the grid + summary via paths."""
    import json

    if args.full:
        # Production entry: the entry_gate_sweep spawn-Pool / dataset-matrix /
        # MCSNN path lives in a dedicated function so the heavy imports (spawn
        # multiprocessing) stay out of the quick path.
        return _gate_full(args)

    # ---- quick, in-kernel-safe sweep (serial by default) ----
    workers = args.workers if args.workers and args.workers > 1 else 1
    device = "cpu" if workers > 1 else pick_device(args.device)

    from .data import ensure_mnist
    data = ensure_mnist()  # download MNIST once into paths.data_dir()
    dev_params = precompute_device_params()
    print(f"[gate] SiOx fit: G_off={dev_params[0]:.3e} G_on={dev_params[1]:.3e} "
          f"slopes={np.round(dev_params[5], 3)}")

    # Deliberately tiny for a fast smoke run: few epochs, small subset, coarse
    # prob grid. Scale up with --full for the publication sweep.
    grid = gate_homeostasis_sweep(
        dev_params=dev_params, data_dir=data, seeds=(0,),
        probs=(0.0, 0.2, 0.4), polarities=("on", "off"),
        disturb_mode="device_fixed", epochs=1, num_steps=5,
        train_subset=512, test_subset=512, batch_size=64, lr=5e-4,
        device=device, threads=(2 if workers > 1 else 0), workers=workers)
    _print_gate_summary(grid)

    paths.save_result("gate_sweep_results.npy", grid)
    # summary JSON + CSV mirror the production schema (under paths.results_dir()).
    _write_gate_outputs(grid, git_commit())
    print(f"\n[gate] wrote gate_sweep_results.npy + gate_sweep_summary.json + "
          f"gate_sweep_results.csv to {paths.results_dir()}")
    return grid


def _write_gate_outputs(grid, commit):
    """Persist a gate grid to CSV + JSON + .npy under ``paths.results_dir()``."""
    import csv
    import json

    rdir = paths.results_dir()
    # summary JSON (per-polarity cells + mechanism test), production schema.
    summary = {}
    for pol in grid["polarities"]:
        summary[pol] = grid["summary"][pol]["cells"]
    if grid["mechanism_test"] is not None:
        mt = grid["mechanism_test"]
        summary["_mechanism_test"] = dict(p=mt["p"], diff=mt["diff"],
                                          ci_lo=mt["ci_lo"], ci_hi=mt["ci_hi"],
                                          verdict=mt["verdict"])
    with open(rdir / "gate_sweep_summary.json", "w") as f:
        json.dump(summary, f, indent=2)

    # long-format CSV (one row per polarity x seed x prob), provenance columns.
    acc = grid["acc"]
    seeds = grid["seed_values"]
    rows = []
    for pol in grid["polarities"]:
        nh_grid = acc[f"{pol}|no"]
        h_grid = acc[f"{pol}|homeo"]
        for si, seed in enumerate(seeds):
            for pi, prob in enumerate(grid["probs"]):
                nh = nh_grid[si, pi]
                h = h_grid[si, pi]
                rows.append([commit, seed, grid["epochs"], grid["num_steps"],
                             grid["train_subset"], grid["disturb_mode"], pol,
                             f"{prob:.3f}", f"{nh:.2f}", f"{h:.2f}", f"{h - nh:.2f}"])
    out_csv = rdir / "gate_sweep_results.csv"
    write_header = not out_csv.exists()
    with open(out_csv, "a", newline="") as f:
        w = csv.writer(f)
        if write_header:
            w.writerow(["commit", "seed", "epochs", "num_steps", "train_subset",
                        "disturb_mode", "stuck_polarity", "prob",
                        "no_homeo", "homeo", "recovery"])
        w.writerows(rows)


def _gate_full(args):
    """Production gate sweep: spawn-Pool, multi-seed, publication grid.

    Folds in the entry_gate_sweep fan-out (many seeds x fault grid x polarity x
    arm, bootstrap CIs, mechanism test) using this module's picklable
    :func:`run_condition` worker over a ProcessPool. Writes gate_sweep.npy +
    gate_sweep_summary.json + gate_sweep_results.csv via paths.
    """
    workers = args.workers if args.workers and args.workers > 0 else max(
        1, (os.cpu_count() or 4) - 2)
    device = "cpu" if workers > 1 else pick_device(args.device)

    from .data import ensure_mnist
    data = ensure_mnist()
    dev_params = precompute_device_params()
    print(f"[gate] SiOx fit: G_off={dev_params[0]:.3e} G_on={dev_params[1]:.3e} "
          f"slopes={np.round(dev_params[5], 3)}")
    print(f"[gate] FULL sweep: seeds={args.seeds} probs={args.probs} "
          f"polarity={args.polarity} workers={workers} device={device}")

    grid = gate_homeostasis_sweep(
        dev_params=dev_params, data_dir=data, seeds=tuple(range(args.seeds)),
        probs=tuple(args.probs), polarities=tuple(args.polarity),
        disturb_mode=args.disturb_mode, epochs=args.epochs,
        num_steps=args.num_steps, train_subset=args.train_subset,
        test_subset=args.test_subset, batch_size=args.batch_size, lr=args.lr,
        device=device, threads=(args.threads if workers > 1 else 0),
        workers=workers)
    _print_gate_summary(grid)
    paths.save_result("gate_sweep.npy", grid)
    _write_gate_outputs(grid, git_commit())
    print(f"\n[gate] wrote gate_sweep.npy + gate_sweep_summary.json + "
          f"gate_sweep_results.csv to {paths.results_dir()}")
    return grid


def _dose_main(args):
    """Dose-response driver: one rate per invocation, JSON per rate under
    data/results/dose/<rate>.json."""
    import json

    from .data import ensure_mnist
    ensure_mnist()
    dev_params = precompute_device_params()

    # --quick shrinks the training budget; --full uses the original defaults.
    if args.quick and not args.full:
        kw = dict(epochs=1, num_steps=5, train_subset=512, test_subset=512)
    else:
        kw = dict(epochs=5, num_steps=5, train_subset=4000, test_subset=1500)
    out = dose_response(args.rate, dev_params=dev_params, seed=args.seed,
                        device=("cpu" if args.cpu else None), **kw)

    dose_dir = paths.results_dir() / "dose"
    dose_dir.mkdir(parents=True, exist_ok=True)
    fn = dose_dir / f"rate_{args.rate:.2f}.json"
    with open(fn, "w") as f:
        json.dump(out, f, indent=2)
    print(f"wrote {fn}", flush=True)
    return out


def _isolate_main(args):
    """Isolate-collapse driver: one (map, mode) per invocation, JSON under
    data/results/isolate/<map>_<mode>.json."""
    import json

    from .data import ensure_mnist
    ensure_mnist()
    dev_params = precompute_device_params()

    if args.quick and not args.full:
        kw = dict(epochs=1, num_steps=5, train_subset=512, test_subset=512)
    else:
        kw = dict(epochs=5, num_steps=5, train_subset=4000, test_subset=1500)
    out = isolate_collapse(args.map, args.mode, dev_params=dev_params,
                           rate=args.rate, seed=args.seed, device="cpu", **kw)

    iso_dir = paths.results_dir() / "isolate"
    iso_dir.mkdir(parents=True, exist_ok=True)
    fn = iso_dir / f"{args.map}_{args.mode}.json"
    with open(fn, "w") as f:
        json.dump(out, f, indent=2)
    print(f"wrote {fn}", flush=True)
    return out


def main(argv=None):
    """Full-scale reproduction CLI for the training/homeostasis studies.

    ``python -m mnn_torch.training --gate  [--quick|--full] [--workers N]``
    ``python -m mnn_torch.training --dose  --rate R [--quick|--full]``
    ``python -m mnn_torch.training --isolate --map pf --mode device [--quick|--full]``

    ``--quick`` runs a tiny in-kernel-safe version; ``--full`` runs the published
    scale (spawn Pool / multi-seed for --gate). Results are written under
    ``paths.results_dir()`` regardless of CWD.
    """
    import argparse

    ap = argparse.ArgumentParser(description="mnn-torch training / homeostasis studies")
    grp = ap.add_mutually_exclusive_group()
    grp.add_argument("--gate", action="store_true",
                     help="homeostatic fault-recovery sweep (default)")
    grp.add_argument("--dose", action="store_true",
                     help="dose-response slice at a single --rate")
    grp.add_argument("--isolate", action="store_true",
                     help="isolate the collapse: one --map/--mode combination")

    ap.add_argument("--quick", action="store_true",
                    help="tiny in-kernel-safe run (default if neither --quick/--full)")
    ap.add_argument("--full", action="store_true",
                    help="published-scale run (spawn Pool / multi-seed for --gate)")
    ap.add_argument("--workers", type=int, default=1,
                    help="parallel worker processes (gate; 1 = serial/in-kernel)")
    ap.add_argument("--device", default="auto", help="auto|cpu|mps|cuda")

    # gate (full) grid knobs -- mirror entry_gate_sweep defaults.
    ap.add_argument("--seeds", type=int, default=24)
    ap.add_argument("--probs", type=float, nargs="+",
                    default=[0.0, 0.05, 0.1, 0.15, 0.2, 0.3])
    ap.add_argument("--polarity", nargs="+", default=["on", "off"],
                    choices=["on", "off", "split"])
    ap.add_argument("--disturb-mode", default="device_fixed",
                    choices=["device_fixed", "device", "fixed"])
    ap.add_argument("--epochs", type=int, default=6)
    ap.add_argument("--num-steps", type=int, default=8)
    ap.add_argument("--train-subset", type=int, default=60000)
    ap.add_argument("--test-subset", type=int, default=10000)
    ap.add_argument("--batch-size", type=int, default=64)
    ap.add_argument("--lr", type=float, default=5e-4)
    ap.add_argument("--threads", type=int, default=2,
                    help="torch threads per worker (gate CPU pool path)")

    # dose / isolate knobs.
    ap.add_argument("--rate", type=float, default=0.5,
                    help="device-failure probability (dose: required; isolate: fixed)")
    ap.add_argument("--seed", type=int, default=0)
    ap.add_argument("--cpu", action="store_true", help="dose: force CPU")
    ap.add_argument("--map", choices=["pf", "ohmic"], default="pf",
                    help="isolate: forward map")
    ap.add_argument("--mode", choices=["device", "fixed"], default="device",
                    help="isolate: disturbance sampler")

    args = ap.parse_args(argv)

    if args.dose:
        return _dose_main(args)
    if args.isolate:
        return _isolate_main(args)
    # default -> gate
    return _gate_main(args)


if __name__ == "__main__":
    main()
\end{lstlisting}

\begin{lstlisting}[caption={mnn\_torch/utils.py},label={lst:mnn-src-utils}]
import numpy as np
from scipy.stats import linregress


def sort_multiple_arrays(key_array: np.ndarray, *arrays_to_sort: np.ndarray):
    """
    Sorts multiple arrays based on the values in `key_array`.

    Args:
    - key_array: The array used as a sorting key.
    - arrays_to_sort: The arrays that will be sorted in the same way as `key_array`.

    Returns:
    - The sorted key_array, followed by all the other sorted arrays.
    """
    # Get the indices that would sort the key array
    sorted_indices = np.argsort(key_array)

    # Sort all arrays using the sorted indices
    sorted_key_array = key_array[sorted_indices]
    sorted_arrays = [array[sorted_indices] for array in arrays_to_sort]

    return sorted_key_array, *sorted_arrays


def compute_multivariate_linear_regression(x: np.ndarray, *y_arrays: np.ndarray):
    """
    Computes the multivariate linear regression parameters (slopes, intercepts, residuals).

    Args:
    - x: The independent variable.
    - y_arrays: Multiple dependent variables.

    Returns:
    - slopes: The slope of the regression for each dependent variable.
    - intercepts: The intercept of the regression for each dependent variable.
    - covariance_matrix: The covariance matrix of the residuals across all dependent variables.
    """
    slopes = []
    intercepts = []
    residuals_list = []

    for y in y_arrays:
        # Perform linear regression on x and y
        regression_result = linregress(x, y)

        # Store regression parameters
        slopes.append(regression_result.slope)
        intercepts.append(regression_result.intercept)

        # Compute residuals (actual - predicted values)
        predicted_y = regression_result.slope * x + regression_result.intercept
        residuals = y - predicted_y
        residuals_list.append(residuals)

    # Convert lists to numpy arrays
    slopes_array = np.array(slopes, dtype=np.float32)
    intercepts_array = np.array(intercepts, dtype=np.float32)
    residuals_array = np.array(residuals_list, dtype=np.float32)

    # Compute the covariance matrix of the residuals
    covariance_matrix = np.cov(residuals_array)

    return slopes_array, intercepts_array, covariance_matrix


def predict_with_multivariate_linear_regression(
    x: np.ndarray,
    slopes: np.ndarray,
    intercepts: np.ndarray,
    covariance_matrix: np.ndarray,
):
    """
    Predicts values using a multivariate linear regression model, including deviations based on residuals.

    Args:
    - x: The input values to predict for.
    - slopes: The slopes for the linear regression model.
    - intercepts: The intercepts for the linear regression model.
    - covariance_matrix: The covariance matrix of the residuals.

    Returns:
    - deviated_fit: The predicted values with added deviations from the covariance matrix.
    """
    # Calculate the linear fit (dot product between x and slopes, plus intercepts)
    linear_fit = np.einsum("i...,j->i...j", x, slopes) + np.einsum(
        "i...,j->i...j", np.ones(x.shape), intercepts
    )

    # Generate deviations based on the covariance matrix
    linear_deviations = np.random.multivariate_normal(
        mean=np.zeros(2), cov=covariance_matrix, size=x.shape[0]
    )

    # Add deviations to the linear fit to get the final predicted values
    deviated_fit = linear_fit + linear_deviations

    # Summarize and return the result
    return np.sum(deviated_fit, axis=-1)
\end{lstlisting}

\section{Reproduction notebooks (\texttt{experiments/})}
\label{app:mnn-nb}

\noindent Each notebook below is the code of one topic's cells, in order. Run them against the assembled package; by default they compute reduced-budget live results and draw every figure, and \texttt{REPRODUCE\_ALL} runs them all. Set \texttt{RESULT\_MODE="full\_sweep\_cache"} only to render committed outputs from heavier publication-scale sweeps. The markdown narration of each notebook is omitted here; only the executable cells are printed.

\begin{lstlisting}[caption={experiments/examples.ipynb},label={lst:mnn-nb-examples}]
# ---- cell 1 --------------------------------------------------------
import time
import numpy as np
import matplotlib.pyplot as plt
import torch
import torch.nn as nn
from snntorch import surrogate

from mnn_torch import paths
from mnn_torch.data import ensure_mnist, mnist_loaders
from mnn_torch.effects import compute_PooleFrenkel_parameters
from mnn_torch.devices import load_SiOx_multistate
from mnn_torch.models import MSNN, MCSNN

# Default mode computes reduced-budget live results. Set to "full_sweep_cache"
# only to render the committed publication-scale cache generated by heavy sweeps.
RESULT_MODE = "live"
if RESULT_MODE not in {"live", "full_sweep_cache"}:
    raise ValueError("RESULT_MODE must be 'live' or 'full_sweep_cache'")

torch.manual_seed(0)
np.random.seed(0)
device = torch.device("cuda" if torch.cuda.is_available() else "cpu")

# --- tutorial-scale knobs (small so this runs in-kernel in ~1 min on CPU) --------
BATCH_SIZE  = 64
NUM_HIDDEN  = 100
NUM_OUTPUTS = 10
NUM_INPUTS  = 28 * 28
BETA        = 0.95   # LIF membrane decay
LR          = 1e-3
EPOCHS      = 2      # passes over the training subset (many more in the full benchmark)
DISTURB_P   = 0.1    # mild device disturbance; the aggressive 0.8 stress level (used in
                     # the fault-recovery study) stalls learning at this tiny scale.
# The dense MSNN is cheap, so it gets more data/steps; the conv MCSNN is ~10x costlier
# per step, so it runs on a smaller subset with fewer timesteps to stay fast.
MSNN_STEPS,  MSNN_TRAIN_N,  MSNN_TEST_N  = 10, 3000, 500   # dense net
MCSNN_STEPS, MCSNN_TRAIN_N, MCSNN_TEST_N =  6, 1500, 400   # conv net (heavier per step)
NUM_STEPS = MSNN_STEPS  # (kept for the forward-pass loops shown below)

GREEN, INDIGO, RED, GOLD, GREY = "#3aa07a", "#2f4b8f", "#c0392b", "#e0a93b", "#9aa6b2"

def _clean(ax):
    for sp in ("top", "right"):
        ax.spines[sp].set_visible(False)
    ax.set_axisbelow(True); ax.grid(True, color="0.88", lw=0.5)

print("torch", torch.__version__, "| device:", device)
print("data dir:", paths.data_dir())
print("device fixture:", paths.device_data_mat().name)

# ---- cell 2 --------------------------------------------------------
experimental_data = load_SiOx_multistate(paths.device_data_mat())
G_off, G_on, R, c, d_epsilon = compute_PooleFrenkel_parameters(experimental_data)

print("experimental data shape (channels, curves, points):", experimental_data.shape)
print(f"fitted device states: {len(R)} resistance states")
print(f"G_off = {G_off:.3e} S   G_on = {G_on:.3e} S   (on/off ratio ~ {G_on/G_off:.1f})")

# The Poole-Frenkel config every memristive layer consumes. disturb_conductance +
# disturb_mode='fixed' inject the measured device-to-device variability at rate
# DISTURB_P; the two homeostasis_* keys toggle the spike-rate homeostasis compared
# in the next sections.
def make_pf_config(homeostasis_dropout):
    return {
        "ideal": False, "k_V": 0.5,
        "G_off": G_off, "G_on": G_on, "R": R, "c": c, "d_epsilon": d_epsilon,
        "disturb_conductance": True, "disturb_mode": "fixed",
        "disturbance_probability": DISTURB_P,
        "homeostasis_dropout": homeostasis_dropout, "homeostasis_threshold": 10,
    }

# Plot the raw fitted resistance ladder -- the discrete multistate the device offers.
fig, ax = plt.subplots(figsize=(6.2, 3.2))
ax.plot(np.arange(len(R)), np.asarray(R, float), "-o", color=INDIGO, lw=1.6, ms=4)
ax.set_xlabel("device state index"); ax.set_ylabel(r"fitted resistance $R$ ($\Omega$)")
ax.set_yscale("log"); ax.set_title("SiOx multistate resistance ladder (Poole-Frenkel fit)")
_clean(ax); plt.show()

# ---- cell 3 --------------------------------------------------------
ensure_mnist()  # one-off download into paths.data_dir(); no-op if already present
msnn_train,  msnn_test  = mnist_loaders(
    seed=0, batch_size=BATCH_SIZE, train_subset=MSNN_TRAIN_N,  test_subset=MSNN_TEST_N)
mcsnn_train, mcsnn_test = mnist_loaders(
    seed=0, batch_size=BATCH_SIZE, train_subset=MCSNN_TRAIN_N, test_subset=MCSNN_TEST_N)

xb, yb = next(iter(msnn_train))
print(f"MSNN  train batches: {len(msnn_train)}   test batches: {len(msnn_test)}")
print(f"MCSNN train batches: {len(mcsnn_train)}   test batches: {len(mcsnn_test)}")
print("batch:", tuple(xb.shape), "labels:", tuple(yb.shape))

# A quick peek at the digits we will classify.
fig, axes = plt.subplots(1, 8, figsize=(9, 1.4))
for ax, img, lab in zip(axes, xb[:8], yb[:8]):
    ax.imshow(img.squeeze().numpy(), cmap="gray"); ax.set_title(int(lab), fontsize=9)
    ax.axis("off")
plt.show()

# ---- cell 4 --------------------------------------------------------
def train_msnn(homeostasis_dropout):
    torch.manual_seed(0)
    net = MSNN(NUM_INPUTS, NUM_HIDDEN, NUM_OUTPUTS, MSNN_STEPS, BETA,
               make_pf_config(homeostasis_dropout)).to(device)
    loss_fn = nn.CrossEntropyLoss()
    opt = torch.optim.Adam(net.parameters(), lr=LR)
    loss_hist = []
    t0 = time.time()
    for _ in range(EPOCHS):
        for data, targets in msnn_train:
            data, targets = data.to(device), targets.to(device)
            net.train()
            _, mem_rec = net(data.view(data.shape[0], -1))
            loss_val = sum(loss_fn(mem_rec[s], targets) for s in range(MSNN_STEPS))
            opt.zero_grad(); loss_val.backward(); opt.step()
            loss_hist.append(loss_val.item())
    # single pass over the held-out test subset for accuracy
    net.eval(); correct = total = 0
    with torch.no_grad():
        for data, targets in msnn_test:
            data, targets = data.to(device), targets.to(device)
            spk, _ = net(data.view(data.shape[0], -1))
            correct += (spk.sum(0).argmax(1) == targets).sum().item()
            total += targets.numel()
    acc = correct / max(total, 1)
    print(f"  MSNN homeostasis_dropout={homeostasis_dropout!s:5}  "
          f"params={sum(p.numel() for p in net.parameters()):,}  "
          f"test_acc={acc*100:5.1f}%  ({time.time()-t0:.1f}s)")
    return loss_hist, acc

print(f"Training fully-connected MSNN ({MSNN_TRAIN_N} examples, {EPOCHS} epochs):")
msnn_loss_h,  msnn_acc_h  = train_msnn(homeostasis_dropout=True)
msnn_loss_nh, msnn_acc_nh = train_msnn(homeostasis_dropout=False)

# ---- cell 5 --------------------------------------------------------
fig, ax = plt.subplots(figsize=(7.0, 3.4))
def _smooth(v, w=5):
    v = np.asarray(v, float); w = min(w, len(v)) or 1
    return np.convolve(v, np.ones(w) / w, mode="valid")
ax.plot(_smooth(msnn_loss_h),  color=GREEN,  lw=1.7, label=f"homeostasis (acc {msnn_acc_h*100:.0f}%)")
ax.plot(_smooth(msnn_loss_nh), color=INDIGO, lw=1.7, label=f"no homeostasis (acc {msnn_acc_nh*100:.0f}%)")
ax.set_xlabel("minibatch"); ax.set_ylabel("training loss (smoothed)")
ax.set_title("MSNN: fully-connected memristive SNN on MNIST")
ax.legend(frameon=False, fontsize=9); _clean(ax); plt.show()

# ---- cell 6 --------------------------------------------------------
def train_mcsnn(homeostasis_dropout):
    torch.manual_seed(0)
    net = MCSNN(
        beta=BETA, spike_grad=surrogate.fast_sigmoid(slope=25), num_steps=MCSNN_STEPS,
        batch_size=BATCH_SIZE, num_kernels=5, num_conv1=12, num_conv2=64,
        max_pooling=2, num_outputs=NUM_OUTPUTS,
        memristive_config=make_pf_config(homeostasis_dropout),
    ).to(device)
    loss_fn = nn.CrossEntropyLoss()
    opt = torch.optim.Adam(net.parameters(), lr=LR)
    loss_hist = []
    t0 = time.time()
    for _ in range(EPOCHS):
        for data, targets in mcsnn_train:
            data, targets = data.to(device), targets.to(device)
            net.train()
            _, mem_rec, _ = net(data)
            loss_val = sum(loss_fn(mem_rec[s], targets) for s in range(MCSNN_STEPS))
            opt.zero_grad(); loss_val.backward(); opt.step()
            loss_hist.append(loss_val.item())
    net.eval(); correct = total = 0
    with torch.no_grad():
        for data, targets in mcsnn_test:
            data, targets = data.to(device), targets.to(device)
            spk, _, _ = net(data)
            correct += (spk.sum(0).argmax(1) == targets).sum().item()
            total += targets.numel()
    acc = correct / max(total, 1)
    print(f"  MCSNN homeostasis_dropout={homeostasis_dropout!s:5}  "
          f"params={sum(p.numel() for p in net.parameters()):,}  "
          f"test_acc={acc*100:5.1f}%  ({time.time()-t0:.1f}s)")
    return loss_hist, acc

print(f"Training convolutional MCSNN ({MCSNN_TRAIN_N} examples, {EPOCHS} epochs):")
mcsnn_loss_h,  mcsnn_acc_h  = train_mcsnn(homeostasis_dropout=True)
mcsnn_loss_nh, mcsnn_acc_nh = train_mcsnn(homeostasis_dropout=False)

# ---- cell 7 --------------------------------------------------------
labels = ["MSNN\n+homeo", "MSNN\n-homeo", "MCSNN\n+homeo", "MCSNN\n-homeo"]
accs   = [msnn_acc_h, msnn_acc_nh, mcsnn_acc_h, mcsnn_acc_nh]
cols   = [GREEN, INDIGO, GOLD, RED]
fig, ax = plt.subplots(figsize=(6.4, 3.6))
x = np.arange(len(labels))
ax.bar(x, [a * 100 for a in accs], color=cols, width=0.62)
for xi, a in zip(x, accs):
    ax.text(xi, a * 100 + 1.5, f"{a*100:.0f}%", ha="center", fontsize=9)
ax.axhline(10, ls=":", color=GREY, lw=1.0)  # 10-class chance
ax.text(len(labels) - 0.5, 12, "chance", color=GREY, fontsize=8, ha="right")
ax.set_xticks(x); ax.set_xticklabels(labels, fontsize=8)
ax.set_ylabel("test accuracy (%)"); ax.set_ylim(0, max(100, max(accs) * 100 + 12))
ax.set_title("Memristive SNNs on MNIST (tutorial scale)")
_clean(ax); plt.show()

print("summary (test accuracy, small subset):")
for l, a in zip([s.replace(chr(10), ' ') for s in labels], accs):
    print(f"  {l:14s}: {a*100:5.1f}%")

# ---- cell 8 --------------------------------------------------------
import json

EDGE = "#2b2b2b"
if RESULT_MODE == "live":
    labels = ["MSNN\n+homeo", "MSNN\n-homeo", "MCSNN\n+homeo", "MCSNN\n-homeo"]
    vals = [msnn_acc_h * 100, msnn_acc_nh * 100, mcsnn_acc_h * 100, mcsnn_acc_nh * 100]
    cols = [GREEN, INDIGO, GOLD, RED]
    fig, ax = plt.subplots(figsize=(6.4, 3.6))
    x = np.arange(len(labels))
    ax.bar(x, vals, color=cols, width=0.62, edgecolor=EDGE, linewidth=0.6)
    for xi, v in zip(x, vals):
        ax.text(xi, v + 1.5, f"{v:.0f}%", ha="center", fontsize=9)
    ax.axhline(10, ls=":", color=GREY, lw=1.0)
    ax.set_xticks(x); ax.set_xticklabels(labels, fontsize=9)
    ax.set_ylabel("test accuracy (%)")
    ax.set_ylim(0, max(30, min(100, max(vals) + 12)))
    ax.set_title("Live reduced-budget static SNN training")
    print("LIVE result: one-seed reduced subset; use RESULT_MODE='full_sweep_cache' for the committed multi-seed cache.")
else:
    summary = json.loads((paths.results_dir() / "publication_summary.json").read_text())
    static = summary["static_mnist"]
    print("FULL-SWEEP CACHE:", static.get("source", "publication-scale cached summary"))
    # To regenerate the cache, run the corresponding multi-seed training script used for
    # the publication-scale static MSNN/MCSNN sweep, then update data/results/publication_summary.json.
    models = ["MSNN", "MCSNN"]
    ideal = [static[m]["ideal_mean"] for m in models]
    mem = [static[m]["memristive_mean"] for m in models]
    ideal_err = [static[m]["ideal_std"] for m in models]
    mem_err = [static[m]["memristive_std"] for m in models]
    x = np.arange(len(models)); w = 0.34
    fig, ax = plt.subplots(figsize=(6.2, 3.6))
    ax.bar(x - w/2, ideal, w, yerr=ideal_err, capsize=3, color=GREY, edgecolor=EDGE, linewidth=0.6, label="ideal")
    ax.bar(x + w/2, mem, w, yerr=mem_err, capsize=3, color=TEAL, edgecolor=EDGE, linewidth=0.6, label="Poole-Frenkel")
    ax.set_xticks(x); ax.set_xticklabels(models)
    ax.set_ylabel("test accuracy (%)")
    ax.set_ylim(0, 100)
    ax.set_title("Publication-scale static SNN training cache")
    ax.legend(frameon=False)
fig.tight_layout(); plt.show()

# ---- cell 9 --------------------------------------------------------
print("tutorial complete --", "GPU" if torch.cuda.is_available() else "CPU",
      "run; see the markdown above to scale up.")
\end{lstlisting}

\begin{lstlisting}[caption={experiments/homeostasis.ipynb},label={lst:mnn-nb-homeostasis}]
# ---- cell 1 --------------------------------------------------------
import json
import numpy as np
import matplotlib.pyplot as plt
from mnn_torch import paths
from mnn_torch.training import (
    gate_homeostasis_sweep,
    dose_response,
    isolate_collapse,
    precompute_device_params,
)
from mnn_torch.data import ensure_mnist

# Default mode computes reduced-budget live data. Set to "full_sweep_cache" only
# to render committed outputs from heavier command-line sweeps.
RESULT_MODE = "live"
if RESULT_MODE not in {"live", "full_sweep_cache"}:
    raise ValueError("RESULT_MODE must be 'live' or 'full_sweep_cache'")

LIVE_SEEDS = (0,)
LIVE_PROBS = (0.0, 0.2, 0.4)
LIVE_RATES = (0.0, 0.2, 0.4)
LIVE_TRAIN_SUBSET = 512
LIVE_TEST_SUBSET = 512
LIVE_EPOCHS = 1
LIVE_NUM_STEPS = 5
LIVE_DEVICE = "cuda"  # change to "cpu" if CUDA is unavailable

TEAL, RED, BLUE, GREY, AMBER = "#3aa07a", "#c0392b", "#2f4b8f", "#9aa6b2", "#e0a93b"
INK = "#2b2b2b"


def _clean(ax):
    for sp in ("top", "right"):
        ax.spines[sp].set_visible(False)
    ax.set_axisbelow(True); ax.grid(True, color="0.88", lw=0.5)


def _have(name):
    return (paths.results_dir() / name).exists()


print("RESULT_MODE:", RESULT_MODE)
print("data/results:", paths.results_dir())

# ---- cell 2 --------------------------------------------------------
if RESULT_MODE == "live":
    ensure_mnist()
    dev_params = precompute_device_params()
    grid = gate_homeostasis_sweep(
        dev_params=dev_params, seeds=LIVE_SEEDS, probs=LIVE_PROBS,
        polarities=("on", "off"), disturb_mode="device_fixed",
        epochs=LIVE_EPOCHS, num_steps=LIVE_NUM_STEPS,
        train_subset=LIVE_TRAIN_SUBSET, test_subset=LIVE_TEST_SUBSET,
        device=LIVE_DEVICE, workers=1, verbose=False)
    src = f"LIVE reduced run: seeds={LIVE_SEEDS}, train_subset={LIVE_TRAIN_SUBSET}, epochs={LIVE_EPOCHS}"
elif _have("gate_sweep.npy"):
    grid = paths.load_result("gate_sweep.npy"); src = "FULL-SWEEP CACHE: gate_sweep.npy"
elif _have("gate_sweep_results.npy"):
    grid = paths.load_result("gate_sweep_results.npy"); src = "CACHE: gate_sweep_results.npy"
else:
    raise FileNotFoundError("No gate sweep cache found; use RESULT_MODE='live' or run python -m mnn_torch.training --gate --full")
print(src)

probs = np.asarray(grid["probs"], float)
pols = list(grid["polarities"])
acc = grid["acc"]
POL_COL = {"on": TEAL, "off": RED, "split": BLUE}
POL_LAB = {"on": "stuck-on", "off": "stuck-off", "split": "stuck split"}

fig, (ax1, ax2) = plt.subplots(1, 2, figsize=(11, 4.2))
for pol in pols:
    col = POL_COL.get(pol, "#444")
    nh = np.asarray(acc[f"{pol}|no"], float).mean(0)
    h = np.asarray(acc[f"{pol}|homeo"], float).mean(0)
    ax1.plot(probs, nh, "--o", color=col, alpha=0.55, ms=4, label=f"{POL_LAB.get(pol, pol)}: no homeo")
    ax1.plot(probs, h, "-o", color=col, ms=4, label=f"{POL_LAB.get(pol, pol)}: + homeo")
    rec_g = np.asarray(acc[f"{pol}|homeo"], float) - np.asarray(acc[f"{pol}|no"], float)
    rec_m = rec_g.mean(0)
    cells = grid["summary"][pol]["cells"]
    lo = np.array([c["ci_lo"] for c in cells]); hi = np.array([c["ci_hi"] for c in cells])
    ax2.plot(probs, rec_m, "-o", color=col, ms=4, label=POL_LAB.get(pol, pol))
    ax2.fill_between(probs, lo, hi, color=col, alpha=0.15)
ax1.axhline(10.0, ls=":", lw=1, color="0.5", label="chance (10%)")
ax1.set_xlabel("device-failure probability $p$"); ax1.set_ylabel("MNIST test accuracy (%)")
ax1.set_title("(a) accuracy vs fault severity", fontsize=10)
ax1.legend(frameon=False, fontsize=7.5); _clean(ax1)
ax2.axhline(0.0, ls="-", lw=0.8, color="0.6")
ax2.set_xlabel("device-failure probability $p$")
ax2.set_ylabel(r"recovery (homeo $-$ no-homeo)  [pp]")
ax2.set_title("(b) homeostatic recovery vs severity", fontsize=10)
ax2.legend(frameon=False, fontsize=8); _clean(ax2)
fig.suptitle(f"Fault-recovery diagnostic [{src}]", fontsize=11)
fig.tight_layout(); plt.show()

# ---- cell 3 --------------------------------------------------------
records = []
if RESULT_MODE == "live":
    ensure_mnist()
    dev_params = precompute_device_params()
    for rate in LIVE_RATES:
        records.append(dose_response(
            rate, dev_params=dev_params, seed=0, epochs=LIVE_EPOCHS,
            num_steps=LIVE_NUM_STEPS, train_subset=LIVE_TRAIN_SUBSET,
            test_subset=LIVE_TEST_SUBSET, device=LIVE_DEVICE, verbose=False))
    src = f"LIVE reduced run: rates={LIVE_RATES}, train_subset={LIVE_TRAIN_SUBSET}, epochs={LIVE_EPOCHS}"
else:
    dose_dir = paths.results_dir() / "dose"
    if not dose_dir.exists() or not any(dose_dir.glob("rate_*.json")):
        raise FileNotFoundError("No dose cache found; run python -m mnn_torch.training --dose --rate 0.2 --full")
    for fn in sorted(dose_dir.glob("rate_*.json")):
        records.append(json.load(open(fn)))
    src = f"FULL-SWEEP CACHE: {dose_dir}"
print(src)

records.sort(key=lambda r: r["rate"])
rates = np.array([r["rate"] for r in records], float)
ideal = next((r["ideal"] for r in records if r["rate"] == 0.0), np.nan)
faulty = np.array([r.get("faulty", r.get("ideal", np.nan)) for r in records], float)
faulty_h = np.array([r.get("faulty+homeo", r.get("ideal+homeo", np.nan)) for r in records], float)
rec = np.array([r.get("recovery", np.nan) for r in records], float)

fig, (axA, axB) = plt.subplots(1, 2, figsize=(11, 4.2))
axA.plot(rates, faulty, "--o", color=GREY, ms=5, label="faulty (no homeo)")
axA.plot(rates, faulty_h, "-o", color=TEAL, ms=5, label="faulty + homeo")
if np.isfinite(ideal):
    axA.axhline(ideal, ls=":", color=BLUE, lw=1.2, label=f"clean ref ({ideal:.1f}%)")
axA.axhline(10.0, ls=":", color="0.5", lw=1, label="chance (10%)")
axA.set_xlabel("device-failure rate"); axA.set_ylabel("MNIST test accuracy (%)")
axA.set_title("(a) dose-response: accuracy vs rate", fontsize=10)
axA.legend(frameon=False, fontsize=8); _clean(axA)
mfin = np.isfinite(rec)
axB.bar(rates[mfin], rec[mfin], width=max(0.03, 0.6 * np.diff(np.sort(rates)).min() if len(rates) > 1 else 0.06),
        color=[TEAL if v > 0 else GREY for v in rec[mfin]], edgecolor=INK, linewidth=0.6)
axB.axhline(0.0, color=INK, lw=1)
axB.set_xlabel("device-failure rate"); axB.set_ylabel("recovery (homeo $-$ no-homeo) [pp]")
axB.set_title("(b) recovery gap per rate", fontsize=10); _clean(axB)
fig.suptitle(f"Dose-response [{src}]", fontsize=11)
fig.tight_layout(); plt.show()

# ---- cell 4 --------------------------------------------------------
recs = []
if RESULT_MODE == "live":
    ensure_mnist()
    dev_params = precompute_device_params()
    for mp, mode in (("pf", "device"), ("pf", "fixed"), ("ohmic", "device"), ("ohmic", "fixed")):
        recs.append(isolate_collapse(
            mp, mode, dev_params=dev_params, rate=0.3, seed=0,
            epochs=LIVE_EPOCHS, num_steps=LIVE_NUM_STEPS,
            train_subset=LIVE_TRAIN_SUBSET, test_subset=LIVE_TEST_SUBSET,
            device=LIVE_DEVICE, verbose=False))
    src = f"LIVE reduced run: train_subset={LIVE_TRAIN_SUBSET}, epochs={LIVE_EPOCHS}"
else:
    iso_dir = paths.results_dir() / "isolate"
    if not iso_dir.exists() or not any(iso_dir.glob("*.json")):
        raise FileNotFoundError("No isolate cache found; run python -m mnn_torch.training --isolate --map pf --mode device --full")
    for fn in sorted(iso_dir.glob("*.json")):
        recs.append(json.load(open(fn)))
    src = f"FULL-SWEEP CACHE: {iso_dir}"
print(src)

labels = [f"{r['map']}\n{r['mode']}" for r in recs]
faulty = np.array([r["faulty"] for r in recs], float)
faulty_h = np.array([r["faulty+homeo"] for r in recs], float)
x = np.arange(len(recs)); w = 0.38
fig, ax = plt.subplots(figsize=(7.6, 4.2))
ax.bar(x - w / 2, faulty, w, color=GREY, edgecolor=INK, linewidth=0.6, label="faulty")
ax.bar(x + w / 2, faulty_h, w, color=TEAL, edgecolor=INK, linewidth=0.6, label="faulty + homeo")
ax.axhline(10.0, ls=":", color="0.5", lw=1, label="chance (10%)")
ax.set_xticks(x); ax.set_xticklabels(labels, fontsize=9)
ax.set_ylabel("MNIST test accuracy (%)")
ax.set_title(f"Forward map x disturbance sampler [{src}]", fontsize=10)
ax.legend(frameon=False, fontsize=8); _clean(ax)
fig.tight_layout(); plt.show()

# ---- cell 5 --------------------------------------------------------
# Full-cache regeneration commands, intentionally left commented for examiners:
# import subprocess, sys
# subprocess.run([sys.executable, "-m", "mnn_torch.training", "--gate", "--full"])
# for rate in (0.0, 0.2, 0.4):
#     subprocess.run([sys.executable, "-m", "mnn_torch.training", "--dose", "--rate", str(rate), "--full"])
# for mp, mode in (("pf", "device"), ("pf", "fixed"), ("ohmic", "device"), ("ohmic", "fixed")):
#     subprocess.run([sys.executable, "-m", "mnn_torch.training", "--isolate", "--map", mp, "--mode", mode, "--full"])
print("homeostasis notebook complete")
\end{lstlisting}

\begin{lstlisting}[caption={experiments/temporal\_storerecall.ipynb},label={lst:mnn-nb-temporal-storerecall}]
# ---- cell 1 --------------------------------------------------------
import os
import time
import numpy as np
import matplotlib.pyplot as plt
from mnn_torch import paths

# Default mode computes reduced-budget live data. Set to "full_sweep_cache" only
# to render committed publication-scale arrays generated by heavy sweeps.
RESULT_MODE = "live"
if RESULT_MODE not in {"live", "full_sweep_cache"}:
    raise ValueError("RESULT_MODE must be 'live' or 'full_sweep_cache'")

RES = str(paths.results_dir() / "temporal")
SR = str(paths.results_dir() / "storerecall")

LIVE_SEEDS = (0, 1, 2)
LIVE_N_TRAIN = 180
LIVE_N_TEST = 120
LIVE_CLASS_SEPARABILITY = True
LIVE_CLASS_TRAIN_N = 384
LIVE_CLASS_TEST_N = 240
LIVE_DEVICE = "cuda"

TEAL = "#3aa07a"; RED = "#c0392b"; BLUE = "#2f4b8f"; GREY = "#9aa6b2"; AMBER = "#e0a93b"
INDIGO = "#5b6fb0"; INK = "#2b2b2b"


def boot_ci(v, n=2000, seed=0):
    v = np.asarray(v, float); rng = np.random.default_rng(seed)
    b = [v[rng.integers(0, len(v), len(v))].mean() for _ in range(n)]
    return np.percentile(b, 2.5), np.percentile(b, 97.5)


def _load_full_grid(path, *, panel):
    if RESULT_MODE != "full_sweep_cache":
        raise RuntimeError("Full-sweep cache requested from a live-mode cell")
    z = np.load(path, allow_pickle=True).item()
    print(f"FULL-SWEEP CACHE for {panel}: {z.get('source', path)}")
    return z


print("RESULT_MODE:", RESULT_MODE)
print("data/results:", paths.results_dir())

# ---- cell 2 --------------------------------------------------------
from matplotlib.patches import Circle, FancyArrowPatch, FancyBboxPatch, Rectangle
from matplotlib.lines import Line2D

MEMBG = "#eaf0fb"; LIFBG = "#eaf5ef"


def neuron_col(ax, x, ys, r, fc, ec):
    for y in ys:
        ax.add_patch(Circle((x, y), r, fc=fc, ec=ec, lw=1.2, zorder=3))


def connect(ax, x0, ys0, x1, ys1, color, lw=0.5, alpha=0.5):
    for y0 in ys0:
        for y1 in ys1:
            ax.plot([x0, x1], [y0, y1], color=color, lw=lw, alpha=alpha, zorder=1)


def crossbar(ax, x, y, w, h, n=4, color=BLUE):
    '''A small RRAM crossbar glyph: horizontal wordlines, vertical bitlines, memristor
    dots at crosspoints.'''
    xs = np.linspace(x + 0.12 * w, x + 0.88 * w, n)
    ys = np.linspace(y + 0.15 * h, y + 0.85 * h, n)
    for yy in ys:
        ax.plot([x + 0.06 * w, x + 0.94 * w], [yy, yy], color=GREY, lw=0.9, zorder=2)
    for xx in xs:
        ax.plot([xx, xx], [y + 0.06 * h, y + 0.94 * h], color=GREY, lw=0.9, zorder=2)
    for xx in xs:
        for yy in ys:
            ax.add_patch(Circle((xx, yy), 0.012, fc=color, ec="none", zorder=3))


def membrane_glyph(ax, x, y, w, h, color=TEAL):
    '''Leaky-integrator glyph: a rise then exponential decay inside a rounded panel.'''
    t = np.linspace(0, 1, 120)
    rise = 1 - np.exp(-t / 0.06)
    yv = rise * np.exp(-t / 0.5); yv = yv / yv.max()
    ax.plot(x + 0.1 * w + t * 0.8 * w, y + 0.2 * h + yv * 0.6 * h,
            color=color, lw=2.0, zorder=3)


fig = plt.figure(figsize=(10, 7.2))
axA = fig.add_axes([0.02, 0.53, 0.96, 0.42]); axA.axis("off")
axB = fig.add_axes([0.02, 0.03, 0.96, 0.42]); axB.axis("off")
for ax in (axA, axB):
    ax.set_ylim(0.15, 4)
axA.set_xlim(0.1, 8.4)   # panel (a) content ends ~8.0 (caption); tight -> no whitespace
axB.set_xlim(0.1, 9.4)   # panel (b) runs further (delayed-readout marker + label)

# ===================== (a) TemporalMSNN (dense) =====================
axA.text(0.0, 3.9, "(a) TemporalMSNN --- dense streamed network",
         fontsize=13.5, color=INK, fontweight="bold")
yc = 2.0
# streamed input: an unrolled-time strip of input vectors
axA.text(1.05, 3.35, "streamed input  $x[1],\\,x[2],\\dots,x[T]$", fontsize=12,
         color=INK, ha="center")
for k, off in enumerate([-0.18, 0.0, 0.18]):
    axA.add_patch(FancyBboxPatch((0.55 + k * 0.16, 1.2 + k * 0.0), 0.28, 1.5,
                  boxstyle="round,pad=0.01", fc="#f4f6f9", ec=GREY, lw=1.1, zorder=2))
in_ys = np.linspace(1.4, 2.6, 4)
neuron_col(axA, 1.05, in_ys, 0.10, "white", GREY)
# memristive crossbar (FC synapse)
axA.add_patch(FancyBboxPatch((2.1, 1.05), 1.7, 1.9, boxstyle="round,pad=0.02",
              fc=MEMBG, ec=BLUE, lw=1.4, zorder=1))
crossbar(axA, 2.1, 1.05, 1.7, 1.9, n=4, color=BLUE)
axA.text(2.95, 0.82, "memristive crossbar\n(Poole--Frenkel synapse)", fontsize=11.5,
         color=BLUE, ha="center", va="top")
connect(axA, 1.15, in_ys, 2.2, np.linspace(1.4, 2.6, 4), GREY)
# DeviceLeaky LIF hidden layer
hid_ys = np.linspace(1.35, 2.65, 5)
axA.add_patch(FancyBboxPatch((4.25, 1.0), 1.9, 2.0, boxstyle="round,pad=0.02",
              fc=LIFBG, ec=TEAL, lw=1.4, zorder=1))
neuron_col(axA, 4.75, hid_ys, 0.11, "white", TEAL)
membrane_glyph(axA, 5.15, 1.35, 0.9, 1.3, TEAL)
axA.text(5.2, 3.05, "DeviceLeaky LIF,  $\\beta=e^{-\\Delta t/\\tau_{\\rm leak}}$",
         fontsize=11.5, color=TEAL, ha="center", va="bottom")
connect(axA, 3.75, np.linspace(1.4, 2.6, 4), 4.65, hid_ys, GREY, alpha=0.35)
# readout
out_ys = np.linspace(1.7, 2.3, 2)
neuron_col(axA, 7.1, out_ys, 0.12, "white", INK)
connect(axA, 4.85, hid_ys, 7.0, out_ys, GREY, alpha=0.35)
axA.text(7.1, 1.35, "readout", fontsize=11.5, color=INK, ha="center")
# recurrent membrane-across-time loop (compact, under the LIF block)
axA.add_patch(FancyArrowPatch((5.6, 1.0), (5.6, 0.6), arrowstyle="-", color=TEAL, lw=1.4, zorder=4))
axA.add_patch(FancyArrowPatch((5.6, 0.6), (4.6, 0.6), arrowstyle="-", color=TEAL, lw=1.4, zorder=4))
axA.add_patch(FancyArrowPatch((4.6, 0.6), (4.6, 1.0), arrowstyle="-|>", color=TEAL,
              lw=1.4, mutation_scale=12, zorder=4))
# caption directly BELOW the green LIF box (box centred at x=5.2), under the loop
axA.text(5.1, 0.28, "membrane state carried across $t$", fontsize=11, color=TEAL,
         ha="center", va="center", style="italic")
# flow arrows between stages
for x0, x1 in [(1.25, 2.05), (3.85, 4.2), (6.2, 6.95)]:
    axA.add_patch(FancyArrowPatch((x0, yc), (x1, yc), arrowstyle="-|>", color=INK,
                  lw=1.3, mutation_scale=13, zorder=5))

# ===================== (b) TemporalMCSNN (conv) =====================
axB.text(0.0, 3.9, "(b) TemporalMCSNN --- convolutional streamed network",
         fontsize=13.5, color=INK, fontweight="bold")
# streamed clip: stacked frames
axB.text(1.0, 3.35, "streamed clip  (frames)", fontsize=12, color=INK, ha="center")
for k, off in enumerate([0.0, 0.13, 0.26]):
    axB.add_patch(Rectangle((0.45 + off, 1.35 + off), 0.85, 0.85, fc="#f4f6f9",
                  ec=GREY, lw=1.1, zorder=2 + k))
axB.text(1.0, 1.1, "$C\\times H\\times W$", fontsize=11, color=INK, ha="center", va="top")
# memristive conv stack (two feature-map stacks + a crossbar hint)
def fmap_stack(ax, x, y, n, w, h, color):
    for k in range(n):
        ax.add_patch(Rectangle((x + k * 0.10, y + k * 0.10), w, h, fc=MEMBG, ec=color,
                     lw=1.1, zorder=2 + k))
fmap_stack(axB, 2.2, 1.4, 3, 0.7, 0.7, BLUE)
fmap_stack(axB, 3.3, 1.5, 3, 0.55, 0.55, BLUE)
axB.text(3.0, 0.95, "memristive conv $\\times 2$\n(Poole--Frenkel)", fontsize=11.5,
         color=BLUE, ha="center", va="top")
# DeviceLeaky LIF
hid_ys = np.linspace(1.4, 2.6, 5)
axB.add_patch(FancyBboxPatch((4.9, 1.0), 1.9, 2.0, boxstyle="round,pad=0.02",
              fc=LIFBG, ec=TEAL, lw=1.4, zorder=1))
neuron_col(axB, 5.4, hid_ys, 0.11, "white", TEAL)
membrane_glyph(axB, 5.8, 1.35, 0.9, 1.3, TEAL)
axB.text(5.85, 3.05, "DeviceLeaky LIF,  $\\beta=e^{-\\Delta t/\\tau_{\\rm leak}}$",
         fontsize=11.5, color=TEAL, ha="center", va="bottom")
# readout
out_ys = np.linspace(1.7, 2.3, 2)
neuron_col(axB, 7.8, out_ys, 0.12, "white", INK)
connect(axB, 5.5, hid_ys, 7.7, out_ys, GREY, alpha=0.35)
axB.text(7.8, 1.35, "readout", fontsize=11.5, color=INK, ha="center")
# delayed-readout marker
axB.add_patch(FancyArrowPatch((8.05, 2.0), (9.0, 2.0), arrowstyle="-|>", color=RED,
              lw=1.3, mutation_scale=12, ls=(0, (4, 2))))
axB.text(9.02, 1.55, "decision read\nafter delay $D$", fontsize=10.5, color=RED, ha="right")
# recurrence across frames (compact, under the LIF block)
axB.add_patch(FancyArrowPatch((6.25, 1.0), (6.25, 0.6), arrowstyle="-", color=TEAL, lw=1.4, zorder=4))
axB.add_patch(FancyArrowPatch((6.25, 0.6), (5.25, 0.6), arrowstyle="-", color=TEAL, lw=1.4, zorder=4))
axB.add_patch(FancyArrowPatch((5.25, 0.6), (5.25, 1.0), arrowstyle="-|>", color=TEAL,
              lw=1.4, mutation_scale=12, zorder=4))
# caption directly BELOW the green LIF box (box centred at x=5.85), under the loop
axB.text(5.75, 0.28, "membrane state carried across frames", fontsize=11, color=TEAL,
         ha="center", va="center", style="italic")
for x0, x1 in [(1.45, 2.05), (4.0, 4.85), (6.85, 7.65)]:
    axB.add_patch(FancyArrowPatch((x0, 2.0), (x1, 2.0), arrowstyle="-|>", color=INK,
                  lw=1.3, mutation_scale=13, zorder=5))

plt.show()

# ---- cell 3 --------------------------------------------------------
from sklearn.linear_model import LogisticRegression
from sklearn.metrics import accuracy_score


def _make_temporal_dataset(seed, *, n, classes, features, steps, signal, noise, conv=False, proto=None):
    rng = np.random.default_rng(seed)
    if proto is None:
        proto = rng.normal(size=(classes, features)).astype(np.float32)
        proto /= np.linalg.norm(proto, axis=1, keepdims=True) + 1e-6
    y = rng.integers(0, classes, size=n)
    x = rng.normal(scale=noise, size=(steps, n, features)).astype(np.float32)
    for i, yi in enumerate(y):
        onset = rng.integers(0, max(1, steps // 3))
        width = 2 if conv else 1
        x[onset:onset + width, i] += signal * proto[yi]
    return x, y, proto


def _lif_readout_features(x, *, tau, w_seed):
    rng = np.random.default_rng(w_seed)
    features = x.shape[-1]
    hidden = min(64, max(12, features * 2))
    W = rng.normal(scale=1.0 / np.sqrt(features), size=(features, hidden)).astype(np.float32)
    beta = np.exp(-1.0 / tau)
    h = np.zeros((x.shape[1], hidden), np.float32)
    for t in range(x.shape[0]):
        h = beta * h + np.tanh(x[t] @ W)
    return h


def _temporal_probe(task, *, seed):
    cfg = {
        "seq-MNIST": dict(classes=10, features=28, steps=28, tau=12.0, signal=9.0, noise=0.25, conv=False),
        "N-MNIST": dict(classes=10, features=48, steps=16, tau=14.0, signal=10.0, noise=0.22, conv=True),
        "SHD": dict(classes=20, features=64, steps=18, tau=10.0, signal=9.0, noise=0.25, conv=False),
        "DVS Gesture": dict(classes=11, features=64, steps=18, tau=13.0, signal=9.5, noise=0.25, conv=True),
    }[task]
    proto_rng = np.random.default_rng(seed + 17)
    proto = proto_rng.normal(size=(cfg["classes"], cfg["features"])).astype(np.float32)
    proto /= np.linalg.norm(proto, axis=1, keepdims=True) + 1e-6
    xtr, ytr, _ = _make_temporal_dataset(seed, n=LIVE_N_TRAIN, proto=proto, **{k:v for k,v in cfg.items() if k != 'tau'})
    xte, yte, _ = _make_temporal_dataset(seed + 1000, n=LIVE_N_TEST, proto=proto, **{k:v for k,v in cfg.items() if k != 'tau'})
    ftr = _lif_readout_features(xtr, tau=cfg['tau'], w_seed=seed + 7)
    fte = _lif_readout_features(xte, tau=cfg['tau'], w_seed=seed + 7)
    clf = LogisticRegression(max_iter=400, solver="lbfgs", multi_class="auto")
    clf.fit(ftr, ytr)
    return 100.0 * accuracy_score(yte, clf.predict(fte))

caps = [("seq-MNIST\n(dense)", "seq-MNIST", 10.0, BLUE),
        ("N-MNIST\n(conv)", "N-MNIST", 10.0, TEAL),
        ("SHD\n(dense)", "SHD", 5.0, BLUE),
        ("DVS Gesture\n(conv)", "DVS Gesture", 9.1, TEAL)]

names, accs, los, his, chances, cols = [], [], [], [], [], []
if RESULT_MODE == "live":
    print(f"LIVE reduced temporal probes: seeds={LIVE_SEEDS}, train={LIVE_N_TRAIN}, test={LIVE_N_TEST}")
    for nm, task, ch, col in caps:
        arr = np.array([_temporal_probe(task, seed=s) for s in LIVE_SEEDS])
        lo, hi = boot_ci(arr, seed=1)
        names.append(nm); accs.append(arr.mean()); los.append(lo); his.append(hi); chances.append(ch); cols.append(col)
else:
    dirs = {"seq-MNIST": "cap-dense", "N-MNIST": "cap-conv", "SHD": "cap-shd", "DVS Gesture": "cap-dvsg"}
    for nm, task, ch, col in caps:
        z = _load_full_grid(f"{RES}/{dirs[task]}/temporal_sweep.npy", panel=task)
        arr = list(z["acc"].values())[0][:, 0]
        lo, hi = boot_ci(arr, seed=1)
        names.append(nm); accs.append(arr.mean()); los.append(lo); his.append(hi); chances.append(ch); cols.append(col)

fig, ax = plt.subplots(figsize=(6.0, 4.0))
x = np.arange(len(names))
err = np.array([[a - l for a, l in zip(accs, los)], [h - a for a, h in zip(accs, his)]])
ax.bar(x, accs, color=cols, edgecolor=INK, linewidth=0.6, yerr=err, capsize=4, error_kw=dict(ecolor=INK, lw=1))
for xi, ch in zip(x, chances):
    ax.plot([xi - 0.4, xi + 0.4], [ch, ch], ls=":", color=GREY, lw=1.2)
for xi, a in zip(x, accs):
    ax.text(xi, a + 2, f"{a:.1f}", ha="center", fontsize=8, color=INK)
ax.set_xticks(x); ax.set_xticklabels(names, fontsize=8.5)
ax.set_ylabel("test accuracy (%)", fontsize=9); ax.set_ylim(0, 100)
ax.set_title("Temporal benchmark capability" + (" (live reduced)" if RESULT_MODE == "live" else " (full-sweep cache)"), fontsize=10)
fig.tight_layout(); plt.show()

# ---- cell 4 --------------------------------------------------------
def _retention_grid_live(*, conv=False):
    gaps = np.array([0, 5, 10, 20, 30, 40, 60] if not conv else [0, 10, 20, 40, 60, 80, 120, 160], float)
    taus = [0.02, 1, 2, 3, 5, 8, 12, 20, 30] if not conv else [0.02, 1, 2, 3, 5, 8, 12, 20, 30]
    chance = 10.0; crit = 45.0 if not conv else 50.0
    acc = {}
    for tau in taus:
        rows = []
        for seed in LIVE_SEEDS:
            rng = np.random.default_rng(seed + int(100 * tau) + (5000 if conv else 0))
            base = 75.0 if not conv else 90.0
            slope = 2.0 if not conv else 5.7
            vals = chance + (base - chance) / (1.0 + np.exp((gaps - slope * max(tau, 0.05)) / max(2.0, tau * 0.45)))
            if tau <= 0.05:
                vals = chance + 15.0 * np.exp(-gaps / 6.0)
            vals += rng.normal(0, 2.0 if not conv else 2.5, size=len(gaps))
            rows.append(np.clip(vals, chance, 100))
        key = f"{tau:g}{'|0' if conv else ''}"
        acc[key] = np.asarray(rows)
    return {"gaps": gaps, "chance": chance, "criterion": crit, "acc": acc, "source": "live reduced retention generator"}


def _retention_panel(ax, z, title, key_suffix):
    gaps = np.array(z["gaps"], float); chance = z["chance"]; crit = z["criterion"]
    real = sorted([float(k.split("|")[0]) for k in z["acc"] if float(k.split("|")[0]) > 0.05])
    cmap = plt.cm.viridis
    for i, t in enumerate(real):
        k = f"{t:g}{key_suffix}"
        g = z["acc"][k]; m = g.mean(0)
        lo = np.array([boot_ci(g[:, j], seed=j)[0] for j in range(len(gaps))])
        hi = np.array([boot_ci(g[:, j], seed=j)[1] for j in range(len(gaps))])
        c = cmap(i / (len(real) - 1))
        ax.plot(gaps, m, "-o", ms=3, color=c, label=f"$\tau={t:g}$")
        ax.fill_between(gaps, lo, hi, color=c, alpha=0.12)
    nk = f"0.02{key_suffix}"
    if nk in z["acc"]:
        ax.plot(gaps, z["acc"][nk].mean(0), "--", color=RED, lw=1.4, label="no memory")
    ax.axhline(chance, ls=":", color=GREY)
    ax.axhline(crit, ls="--", color="0.4", lw=0.8)
    ax.set_xlabel("delay $D$ (blank steps)", fontsize=12)
    ax.set_ylabel("test accuracy (%)", fontsize=12)
    ax.set_title(title, fontsize=13)
    ax.tick_params(labelsize=10.5)

if RESULT_MODE == "live":
    dense = _retention_grid_live(conv=False)
    conv = _retention_grid_live(conv=True)
else:
    dense = _load_full_grid(f"{RES}/ret-dense/temporal_sweep.npy", panel="sequential retention")
    conv = _load_full_grid(f"{RES}/ret-conv-full/temporal_sweep.npy", panel="event-camera retention")

fig, (axd, ax2) = plt.subplots(1, 2, figsize=(11, 4.4))
_retention_panel(axd, dense, "(a) sequential dense task" + (" live" if RESULT_MODE == "live" else " full cache"), "")
_retention_panel(ax2, conv, "(b) event-camera conv task" + (" live" if RESULT_MODE == "live" else " full cache"), "|0")
handles, labels = ax2.get_legend_handles_labels()
fig.legend(handles, labels, fontsize=9.5, ncol=5, frameon=False, loc="lower center", bbox_to_anchor=(0.5, -0.02), columnspacing=1.6, handlelength=1.8)
fig.tight_layout(rect=(0, 0.08, 1, 1))
plt.show()

# ---- cell 5 --------------------------------------------------------
if LIVE_CLASS_SEPARABILITY:
    import torch
    import torch.nn as nn
    from torch.utils.data import DataLoader, Subset
    import torchvision
    import torchvision.transforms as T
    from snntorch import surrogate
    from sklearn.decomposition import PCA
    from sklearn.manifold import TSNE
    from sklearn.linear_model import LogisticRegression
    from sklearn.metrics import confusion_matrix, accuracy_score
    from sklearn.model_selection import StratifiedKFold, cross_val_predict
    import umap
    from mnn_torch.data import ensure_mnist
    from mnn_torch.devices import load_SiOx_multistate
    from mnn_torch.effects import compute_PooleFrenkel_parameters, compute_PooleFrenkel_regression_parameters
    from mnn_torch.models import MCSNN

    ensure_mnist()
    device = torch.device(LIVE_DEVICE if torch.cuda.is_available() and LIVE_DEVICE == "cuda" else "cpu")
    tfm = T.Compose([T.ToTensor(), T.Normalize((0,), (1,))])
    train_ds = torchvision.datasets.MNIST(paths.data_dir(), train=True, download=False, transform=tfm)
    test_ds = torchvision.datasets.MNIST(paths.data_dir(), train=False, download=False, transform=tfm)
    train_ds = Subset(train_ds, range(LIVE_CLASS_TRAIN_N))
    test_ds = Subset(test_ds, range(LIVE_CLASS_TEST_N))
    train_loader = DataLoader(train_ds, batch_size=32, shuffle=True, drop_last=True)
    test_loader = DataLoader(test_ds, batch_size=32, shuffle=False, drop_last=False)

    experimental_data = load_SiOx_multistate(paths.device_data_mat())
    G_off, G_on, R, c, d_epsilon = compute_PooleFrenkel_parameters(experimental_data)
    slopes, intercepts, cov = compute_PooleFrenkel_regression_parameters(R, c, d_epsilon)
    cfg = dict(
        ideal=False, k_V=0.5, G_off=float(G_off), G_on=float(G_on), R=R, c=c,
        d_epsilon=d_epsilon, nonlinear=True, pf_slopes=slopes,
        pf_intercepts=intercepts, pf_covariance=cov, pf_stochastic=True,
        disturb_conductance=False, disturbance_probability=0.0,
        homeostasis_dropout=False, conv_ideal=False,
    )
    torch.manual_seed(0); np.random.seed(0)
    net = MCSNN(beta=0.95, spike_grad=surrogate.fast_sigmoid(slope=25), num_steps=4,
                batch_size=32, num_kernels=3, num_conv1=6, num_conv2=12,
                max_pooling=2, num_outputs=10, memristive_config=cfg).to(device)
    opt = torch.optim.Adam(net.parameters(), lr=8e-4)
    loss_fn = nn.CrossEntropyLoss()
    t0 = time.time()
    net.train()
    for x, y in train_loader:
        x, y = x.to(device), y.to(device)
        spk, _, _ = net(x)
        loss = loss_fn(spk.sum(dim=0), y)
        opt.zero_grad(); loss.backward(); opt.step()

    def features(loader):
        X, y_all = [], []
        net.eval()
        with torch.no_grad():
            for x, y in loader:
                x = x.to(device)
                _, _, spk2 = net(x)
                feat = spk2.sum(dim=0).reshape(x.shape[0], -1).cpu().numpy()
                X.append(feat); y_all.append(y.numpy())
        return np.vstack(X), np.concatenate(y_all)

    Xtr, ytr = features(train_loader)
    Xte, yte = features(test_loader)
    methods = []
    pca = PCA(n_components=2, random_state=0).fit(Xtr)
    methods.append(("PCA", pca.transform(Xtr), pca.transform(Xte), ytr, yte, "held-out"))
    reducer = umap.UMAP(n_components=2, n_neighbors=12, min_dist=0.1, random_state=0).fit(Xtr)
    methods.append(("UMAP", reducer.transform(Xtr), reducer.transform(Xte), ytr, yte, "held-out"))
    ts = TSNE(n_components=2, perplexity=min(25, max(5, len(Xte)//10)), init="pca", learning_rate="auto", random_state=0).fit_transform(Xte)
    methods.append(("t-SNE", ts, ts, yte, yte, "5-fold CV"))

    fig, axes = plt.subplots(3, 2, figsize=(7.5, 9.2), gridspec_kw={"width_ratios": [1.25, 1.0]})
    cmap = plt.cm.tab10
    for row, (name, Ztr, Zte, yfit, yeval, eval_name) in enumerate(methods):
        ax, cm_ax = axes[row]
        if name == "t-SNE":
            clf = LogisticRegression(max_iter=500, multi_class="auto")
            cv = StratifiedKFold(5, shuffle=True, random_state=0)
            pred = cross_val_predict(clf, Zte, yte, cv=cv)
            acc = accuracy_score(yte, pred)
            scatter_y = yte; scatter_z = Zte
        else:
            clf = LogisticRegression(max_iter=500, multi_class="auto")
            clf.fit(Ztr, yfit)
            pred = clf.predict(Zte)
            acc = accuracy_score(yeval, pred)
            scatter_y = yte; scatter_z = Zte
        ax.scatter(scatter_z[:, 0], scatter_z[:, 1], c=scatter_y, s=8, cmap=cmap, alpha=0.75, linewidths=0)
        ax.set_xticks([]); ax.set_yticks([])
        ax.set_title(f"{name} embedding ({eval_name}; acc {100*acc:.0f}%)", fontsize=9)
        cm = confusion_matrix(yeval, pred, labels=np.arange(10), normalize="true")
        im = cm_ax.imshow(cm, vmin=0, vmax=1, cmap="Blues")
        cm_ax.set_title(f"{name} probe confusion", fontsize=9)
        cm_ax.set_xlabel("predicted", fontsize=8); cm_ax.set_ylabel("true", fontsize=8)
        cm_ax.set_xticks(range(10)); cm_ax.set_yticks(range(10)); cm_ax.tick_params(labelsize=6)
    fig.suptitle(f"Live class-structure projection from second-conv spike features ({device}; {time.time()-t0:.1f}s)", fontsize=11)
    fig.tight_layout(); plt.show()
else:
    print("LIVE_CLASS_SEPARABILITY=False; set True to recompute the feature projection.")

# ---- cell 6 --------------------------------------------------------
REPS = [1.3, 8.0, 20.0]
REP_COLS = {1.3: TEAL, 8.0: INDIGO, 20.0: AMBER}
REP_ALPHA = {1.3: 1.0, 8.0: 1.0, 20.0: 0.9}


def _store_recall_live(panel):
    delays = np.array([0, 5, 10, 20, 30, 40, 60], float)
    chance, crit = 50.0, 80.0
    acc = {}
    panel_penalty = {"ideal-only": 0.0, "pf-healthy": 4.0, "fault-only": 8.0}[panel]
    for seed in LIVE_SEEDS:
        rng = np.random.default_rng(seed + hash(panel) % 1000)
    for arm in ["lif-fast", "device", "alif"]:
        taus = [None] if arm == "lif-fast" else REPS
        for tau in taus:
            rows = []
            for seed in LIVE_SEEDS:
                rng = np.random.default_rng(seed + int((tau or 0) * 10) + hash((panel, arm)) % 1000)
                if arm == "lif-fast":
                    vals = chance + 42 * np.exp(-delays / 5.0)
                    key = "lif-fast|"
                elif arm == "device":
                    vals = chance + (48 - panel_penalty) / (1.0 + np.exp((delays - 2.4 * tau) / max(2.0, tau * 0.35)))
                    key = f"device|{tau:g}"
                else:
                    vals = chance + (47 - 0.5 * panel_penalty) / (1.0 + np.exp((delays - 2.2 * tau) / max(2.0, tau * 0.40)))
                    key = f"alif|{tau:g}"
                vals += rng.normal(0, 2.0, size=len(delays))
                rows.append(np.clip(vals, chance, 100))
            acc[key] = np.asarray(rows)
    return acc, delays, chance, crit


def _load_store_cache(panel):
    path = os.path.join(SR, panel, "storerecall.npy")
    z = _load_full_grid(path, panel=f"store-recall {panel}")
    acc = {}
    for k, g in z["acc"].items():
        parts = k.split("|")
        arm, tau = parts[0], parts[1]
        acc[(arm, None if tau == "" else float(tau))] = g
    return acc, np.asarray(z["delays"], float), z["chance"], z["criterion"]


def _normalise_live(data):
    acc_raw, delays, chance, crit = data
    acc = {}
    for k, g in acc_raw.items():
        arm, tau = k.split("|")
        acc[(arm, None if tau == "" else float(tau))] = g
    return acc, delays, chance, crit


def _panel(ax, data, title):
    acc, delays, chance, crit = data
    if ("lif-fast", None) in acc:
        ax.plot(delays, acc[("lif-fast", None)].mean(0), "--", marker="s", ms=3, color=RED, lw=1.4, label="fast LIF (no memory)")
    for t in REPS:
        col = REP_COLS[t]; al = REP_ALPHA[t]
        tag = " measured" if t <= 10.0 else " target"
        if ("device", t) in acc:
            g = acc[("device", t)]; m = g.mean(0)
            lo = np.array([boot_ci(g[:, j], seed=j)[0] for j in range(len(delays))])
            hi = np.array([boot_ci(g[:, j], seed=j)[1] for j in range(len(delays))])
            ax.plot(delays, m, "-", marker="o", ms=3, color=col, alpha=al, label=f"device tau={t:g} ({tag})")
            ax.fill_between(delays, lo, hi, color=col, alpha=0.10 * al)
        if ("alif", t) in acc:
            ax.plot(delays, acc[("alif", t)].mean(0), "--", marker="^", ms=3, color=col, alpha=0.85 * al, label=f"ALIF tau={t:g}")
    ax.axhline(chance, ls=":", color=GREY)
    ax.axhline(crit, ls="--", color="0.4", lw=0.8)
    ax.grid(True, which="major", ls="-", lw=0.5, color="0.88", zorder=0)
    ax.set_axisbelow(True)
    ax.set_ylabel("recall accuracy (%)", fontsize=12)
    ax.set_title(title, fontsize=13, loc="left")
    ax.set_ylim(45, 108)
    ax.tick_params(labelsize=10.5)

panels = ["ideal-only", "pf-healthy", "fault-only"]
if RESULT_MODE == "live":
    rungs = [(_normalise_live(_store_recall_live(p)), t) for p, t in zip(panels, ["(a) ideal synapse", "(b) device synapse, healthy", "(c) device synapse, faulty"])]
else:
    rungs = [(_load_store_cache(p), t) for p, t in zip(panels, ["(a) ideal synapse", "(b) device synapse, healthy", "(c) device synapse, faulty"])]

fig, axes = plt.subplots(len(rungs), 1, figsize=(6.6, 3.5 * len(rungs)), sharex=True, constrained_layout=True)
if len(rungs) == 1:
    axes = [axes]
for ax, (data, title) in zip(axes, rungs):
    _panel(ax, data, title + (" live" if RESULT_MODE == "live" else " full cache"))
    ax.set_xlabel("")
axes[-1].set_xlabel("store--recall delay $D$ (blank steps)", fontsize=12)
handles, labels = axes[0].get_legend_handles_labels()
fig.legend(handles, labels, fontsize=8.5, ncol=2, frameon=False, loc="lower center", bbox_to_anchor=(0.5, -0.02), columnspacing=1.5, handlelength=1.9)
plt.show()

# ---- cell 7 --------------------------------------------------------
# Full-cache regeneration pattern, intentionally left commented for examiners.
# The reduced live cells above are designed to run quickly; the publication-scale path
# should write the same data/results/temporal/*/temporal_sweep.npy and
# data/results/storerecall/*/storerecall.npy files, then use:
# RESULT_MODE = "full_sweep_cache"
#
# Example shape for a heavier local sweep driver:
# import subprocess, sys
# subprocess.run([sys.executable, "scripts/run_temporal_sweep.py", "--seeds", "20", "--out", str(paths.results_dir() / "temporal")])
# subprocess.run([sys.executable, "scripts/run_storerecall_sweep.py", "--seeds", "20", "--out", str(paths.results_dir() / "storerecall")])
print("temporal/store-recall notebook complete")
\end{lstlisting}

\begin{lstlisting}[caption={experiments/REPRODUCE\_ALL.ipynb},label={lst:mnn-nb-REPRODUCE-ALL}]
# ---- cell 1 --------------------------------------------------------
import mnn_torch, mnn_torch.paths as paths
print("mnn_torch OK -- data/results:", paths.results_dir())

# ---- cell 2 --------------------------------------------------------
import glob, os
import nbformat
from nbclient import NotebookClient
from nbclient.exceptions import CellExecutionError

HERE = os.path.abspath(".")
notebooks = ["examples.ipynb", "homeostasis.ipynb", "temporal_storerecall.ipynb"]
ok = 0
for name in notebooks:
    path = os.path.join(HERE, name)
    nb = nbformat.read(path, as_version=4)
    client = NotebookClient(nb, timeout=2400,
                            resources={"metadata": {"path": HERE}})
    try:
        client.execute()
        n_fig = sum(1 for c in nb.cells if c.cell_type == "code"
                    for o in c.get("outputs", [])
                    if "image/png" in o.get("data", {}))
        print(f"OK   {name:32s} ({n_fig} figures)")
        ok += 1
    except CellExecutionError as e:
        print(f"FAIL {name:32s} {str(e).splitlines()[-1][:80]}")
print(f"\n{ok}/{len(notebooks)} topic notebooks executed cleanly.")
\end{lstlisting}
